%% Follow comments to support use.

%%%%%%%%%%%%%%%%%%%%%%%%%%%%%%%%%%%%%%%%%%%%%%%%%%%%%%%%%
%% STEP 1: Choose options for MSc / BSc / seminar layout and your bibliographic style
%%%%%%%%%%%%%%%%%%%%%%%%%%%%%%%%%%%%%%%%%%%%%%%%%%%%%%%%%

%%  Language: 
%%      finnish, swedish, or english
%%  Pagination (use twoside by default)  
%%      oneside or twoside,
%%  Study programme / kind of report
%%      csm  = Master's thesis in Computer Science Master's Programme;
%%      tkt = Bachelor's thesis in Computer Science Bachelor's Programme;
%%      seminar = seminar report
%%  For Master's thesis choose your line or track:
%%      (30 cr thesis, 2020 onwards, Master's Programme in Computer Science = csm)
%%      software-track-2020 = Software study track
%%      algorithms-track-2020 = Algorithms study track
%%      networking-track-2020 = Networking study track

\documentclass[finnish,twoside,censored,tkt,]{HYthesisML} 


% If wanted, open new chapters only at right page.
% By default, "openany".
%\PassOptionsToClass{openright,twoside,a4paper}{report}
\PassOptionsToClass{openany,twoside,a4paper}{report}

\usepackage{csquotes}
%%%%%%%%%%%%%%%%%%%%%%%%%%%%%%%%%%%%%%%%%%%%%%%%%%%%%%%%%
%% REFERENCES
%% Some notes on bibliography usage and options:
%% natbib -> you can use, e.g., \citep{} or \parencite{} for (Einstein, 1905); with APA \cite -> Einstein, 1905 without ()
%% maxcitenames=2 -> only 2 author names in text citations, if more -> et al. is used
%% maxbibnames=99 as no great need to suppress the biliography list in a thesis
%% for more information see biblatex package documentation, e.g., from https://ctan.org/pkg/biblatex 

%% Reference style: select one 
%% for APA = Harvard style = authoryear -> (Einstein, 1905) use:
\usepackage[style=authoryear,bibstyle=authoryear,backend=biber,natbib=true,maxnames=99,maxcitenames=2,uniquelist=minyear,giveninits=true,uniquename=mininit]{biblatex}
%% for numeric = Vancouver style -> [1] use:
%\usepackage[style=numeric,bibstyle=numeric,backend=biber,natbib=true,maxbibnames=99,giveninits=true,uniquename=init]{biblatex}
%% for alpahbetic -> [Ein05] use:
%\usepackage[style=alphabetic,bibstyle=alphabetic,backend=biber,natbib=true,maxbibnames=99,giveninits=true,uniquename=init]{biblatex}
%

\addbibresource{bibliography.bib}
% in case you want the final delimiter between authors & -> (Einstein & Zweistein, 1905) 
% \renewcommand{\finalnamedelim}{ \& }
% List the authors in the Bibilipgraphy as Lastname F, Familyname G,
\DeclareNameAlias{sortname}{family-given}
% remove the punctuation between author names in Bibliography 
%\renewcommand{\revsdnamepunct}{ }


%% Block of definitions for fonts and packages for picture management.
%% In some systems, the figure packages may not be happy together.
%% Choose the ones you need.

%\usepackage[utf8]{inputenc} % For UTF8 support, in some systems. Use UTF8 when saving your file.

\usepackage{lmodern}         % Font package, again in some systems.
\usepackage{textcomp}        % Package for special symbols
\usepackage[pdftex]{color, graphicx} % For pdf output and jpg/png graphics
\usepackage{epsfig}
\usepackage{subfigure}
\usepackage[pdftex, plainpages=false]{hyperref} % For hyperlinks and pdf metadata
\usepackage{fancyhdr}        % For nicer page headers
\usepackage{tikz}            % For making vector graphics (hard to learn but powerful)
%\usepackage{wrapfig}        % For nice text-wrapping figures (use at own discretion)
\usepackage{amsmath, amssymb} % For better math
\usepackage{physics}         % Physics notations for quantum mechanics
\singlespacing               %line spacing options; normally use single

\fussy
%\sloppy                      % sloppy and fussy commands can be used to avoid overlong text lines
% if you want to see which lines are too long or have too little stuff, comment out the following lines
% \overfullrule=1mm
% to see more info in the detailed log about under/overfull boxes...
% \showboxbreadth=50 
% \showboxdepth=50



%%%%%%%%%%%%%%%%%%%%%%%%%%%%%%%%%%%%%%%%%%%%%%%%%%%%%%%%%
%% STEP 2:
%%%%%%%%%%%%%%%%%%%%%%%%%%%%%%%%%%%%%%%%%%%%%%%%%%%%%%%%%
%% Set up personal information for the title page and the abstract form.
%% Replace parameters with your information.
\title{Kvanttiturvallinen salaus}

\author{Aleksi Suuronen}
\date{\today}

% Set supervisors, use the titles according to the thesis language
% in English Prof. or Dr., or in Finnish toht. or tri or FT, TkT, Ph.D. or in Swedish... 
\supervisors{Prof.~Nikolaj Tatti}

\keywords{Kvanttiturvallinen salaus, kvanttilaskenta, kvanttitietokoneet, salausjärjestelmät, kryptografia}
\additionalinformation{\translate{\track}}

%% For seminar reports:
%%\additionalinformation{Name of the seminar}

%% Provide classification terms, to appear on the abstract page.
%% Replace the classification terms below with the ones that match your work.
%% ACM Digital library provides a taxonomy and a tool for classification
%% in computer science. Use 1-3 paths, and use right arrows between the
%% about three levels in the path; each path requires a new line.

\classification{\protect{\ \\
\  Security and privacy $\rightarrow$ Cryptography
}}

%% If you want to quote someone special. You can comment this line out and there will be nothing on the document.
%\quoting{Bachelor's degrees make pretty good placemats if you get them laminated.}{Jeph Jacques}


%% OPTIONAL STEP: Set up properties and metadata for the pdf file that pdfLaTeX makes.
%% Your name, work title, and keywords are recommended.
\hypersetup{
    unicode=true,           % to show non-Latin characters in Acrobat’s bookmarks
    pdftoolbar=true,        % show Acrobat’s toolbar?
    pdfmenubar=true,        % show Acrobat’s menu?
    pdffitwindow=false,     % window fit to page when opened
    pdfstartview={FitH},    % fits the width of the page to the window
    pdftitle={Kvanttiturvallinen salaus},            % title
    pdfauthor={Aleksi Suuronen},           % author
    pdfsubject={Bachelor's thesis},          % subject of the document
    pdfcreator={Aleksi Suuronen},          % creator of the document
    pdfproducer={pdfLaTeX}, % producer of the document
    pdfkeywords={something} {something else}, % list of keywords for
    pdfnewwindow=true,      % links in new window
    colorlinks=true,        % false: boxed links; true: colored links
    linkcolor=black,        % color of internal links
    citecolor=black,        % color of links to bibliography
    filecolor=magenta,      % color of file links
    urlcolor=cyan           % color of external links
}

%%-----------------------------------------------------------------------------------

\begin{document}

% Generate title page.
\maketitle

%%%%%%%%%%%%%%%%%%%%%%%%%%%%%%%%%%%%%%%%%%%%%%%%%%%%%%%%%
%% STEP 3:
%%%%%%%%%%%%%%%%%%%%%%%%%%%%%%%%%%%%%%%%%%%%%%%%%%%%%%%%%
%% Write your abstract in the separate file, to be positioned here.
%% You can make several abstract pages (if you want it in different languages),
%% in which case you should also define the language of the abstract,
%% as below.

\begin{abstract}{}

% Tämä dokumentti on tarkoitettu Helsingin yliopiston tietojenkäsittelytieteen osaston opin\-näyt\-teiden ja harjoitustöiden ulkoasun ohjeeksi ja mallipohjaksi. Ohje soveltuu kanditutkielmiin, ohjelmistotuotantoprojekteihin, seminaareihin ja maisterintutkielmiin. Tämän ohjeen lisäksi on seurattava niitä ohjeita, jotka opastavat valitsemaan kuhunkin osioon tieteellisesti kiinnostavaa, syvällisesti pohdittua sisältöä.


%Työn aihe luokitellaan  
% ACM Computing Classification System (CCS) mukaisesti, 
% ks.\ \url{https://dl.acm.org/ccs}. 
% Käytä muutamaa termipolkua (1--3), jotka alkavat juuritermistä ja joissa polun tarkentuvat luokat erotetaan toisistaan oikealle osoittavalla nuolella.

% Tähän tulee aikanaan yhden kappaleen eli n. 100 sanainen tiivistelmä. Tiivistelmässä käydään läpi mitä on tutkittu, miten on tutkittu ja mitä tuloksia on saatu.

Kvanttiturvallisella salauksella tarkoitetaan salausjärjestelmiä, jotka ovat turvallisia klassisilla ja kvanttitietokoneilla tehtyjä hyökkäyksiä vastaan. Kvanttitietokoneiden kehittyessä suuren kokoluokan kvanttitietokoneet uhkaavat nykyisiä salausjärjestelmiä ja erityisesti julkisen avaimen salausta. Kvanttitietokone mahdollistaa eksponentiaalisen laskentatehon klassiseen tietokoneeseen verrattuna, minkä takia julkisen avaimen salauksessa käytettävät laskennalliset ongelmat voidaan ratkaista tehokkaasti. Yksi tällaisista kvanttialgoritmeista on Shorin algoritmi, joka ratkaisee julkisen avaimen salaukseen liittyvät matemaattiset ongelmat kuten suurten kokonaislukujen tekijöihinjaon tehokkaasti.

Tutkielmassa esitellään kolme kvanttiturvallista salausmenetelmää; koodipohjainen, hilapohjainen ja monimuuttujapohjainen salaus. Näistä salausmenetelmistä etsitään kvanttiturvallista korvaajaa nykyiselle julkisen avaimen salaukselle. Lisäksi esitellään nykypäivänä käytettyjä salausmenetelmiä ja niiden kvanttiturvallisuutta ja perehdytään myös kvanttilaskentaan ja Shorin algoritmiin.

\end{abstract}


%\begin{otherlanguage}{english}
%\begin{abstract}
%Write your abstract here.

%In addition, make sure that all the entries in this form are completed.

%Finally, specify 1--3 ACM Computing Classification System (CCS) topics, as per \url{https://dl.acm.org/ccs}.
%Each topic is specified with one path, as shown in the example below, and elements of the path separated with an arrow.
%Emphasis of each element individually can be indicated
%by the use of bold face for high importance or italics for intermediate
%level.

%\end{abstract}
%\end{otherlanguage}


% Place ToC
%\newpage
\mytableofcontents

\mainmatter

%%%%%%%%%%%%%%%%%%%%%%%%%%%%%%%%%%%%%%%%%%%%%%%%%%%%%%%%%
%% STEP 4: Write the thesis.
%%%%%%%%%%%%%%%%%%%%%%%%%%%%%%%%%%%%%%%%%%%%%%%%%%%%%%%%%
%% Your actual text starts here. You shouldn't mess with the code above the line except
%% to change the parameters. Removing the abstract and ToC commands will mess up stuff.
%%
%% Command \include{file} includes the file of name file.tex.
%% A new page will be created at every \include command, 
%% which makes it appropriate to use it for large entities such as book chapters. Cannot be nested.
%% It is useful for a big project, as changing one of the include targets 
%% won't force the regeneration of the outputs of all the rest.
%% Alternatively, \input is a more lower level macro 
%% which simply inputs the content of the given file like it was copy&pasted there manually.

\chapter{Johdanto\label{intro}}
Kvanttitietokoneiden rakentamisessa on edistytty tasaisesti viime vuosina (\cite{alagic2020status}). Suuren kokoluokan kvanttitietokoneiden yleistyminen uhkaa tietoturvaa niiden tuoman suuren laskentatehon takia. Etenkin nykyiset julkisen avaimen salaukseen liittyvät salausjärjestelmät ovat uhattuna muuttua turvattomiksi kvanttitietokoneiden myötä.

Julkisen avaimen salaus on laajalti nykypäivänä käytössä oleva salausjärjestelmä, jolla salataan muun muassa tietoliikennettä. %oma viite, yleisesti tiedossa oleva asia
Monet julkisen avaimen salausjärjestelmät perustuvat matemaattisiin ongelmiin, jotka on helppo laskea yhteen suuntaan, mutta niiden peruuttaminen tai käänteisoperaatio on taas vaikeaa (\cite{mavroeidis2018impact}). Esimerkki tällaisesta ongelmasta on suurten kokonaislukujen tekijöihinjako. Suurten kokonaislukujen tekijöihinjako on vaikea laskennallinen ongelma ja sille ei tunneta tehokkaita eli polynomisessa ajassa ratkeavia ratkaisuja nykyisillä tietokoneilla (\cite{doi:10.1137/S0036144598347011}). Vaikealla tarkoitetaan tässä tapauksessa sitä, että ongelman ratkeamista joudutaan odottamaan niin kauan, että lopputulosta ei välttämättä saada missään äärellisessä ajassa. Esimerkiksi salauksen purkamisessa kestää niin kauan, että tiedolla ei ole enää merkitystä kun laskenta on loppunut. %ehkä viite
RSA on yksi ensimmäisiä ja käytetyimpiä julkisen avaimen salausjärjestelmiä, joka perustuu suurten kokonaislukujen tekijöihinjakoon (\cite{montgomery1994survey}).

Erityisesti Shorin algoritmi uhkaa muuttaa julkisen avaimen salauksen turvattomaksi (\cite{mavroeidis2018impact}). Vuonna 1994 matemaatikko Peter Shor esitteli kvanttialgoritmin, jolla julkisen avaimen salaukseen perustuvat vaikeat matemaattiset ongelmat, kuten suurten kokonaislukujen tekijöihinjako, voidaan ratkaista tehokkaasti eli polynomisessa ajassa kvanttitietokoneella (\cite{doi:10.1137/S0036144598347011}).

Kvanttitietokoneiden tulemisen myötä on alettu tutkia kvanttiturvallista salausta (\emph{post-quantum cryptography}) (\cite{alagic2020status}). Kvanttiturvallisella salauksella tarkoitetaan salausjärjestelmiä, jotka ovat turvallisia klassisilla tietokoneilla sekä kvanttitietokoneilla tehtyjä hyökkäyksiä vastaan ja jotka voidaan ottaa käyttöön nykyisten salausjärjestelmien tilalle ilman suurempia muutoksia esimerkiksi nykyisessä tietoliikenteessä ja eri tietoliikenneprotokollissa.

Erilaisia kvanttiturvallisia salausmenetelmiä on tutkittu useita (\cite{mavroeidis2018impact}) ja vanhimmat ideat ovat peräisin 1970-luvulta (\cite{repka2014overview}). Monet kvanttiturvalliset salausjärjestelmät perustuvat laskennallisesti vaativiin ongelmiin, jotka ovat vaikeita ratkaista klassisilla tietokoneilla että kvanttitietokoneilla. %oma juttu
 Lisäksi nämä laskennalliset ongelmat ovat luonteeltaan sellaisia, että niihin ei ole osattu soveltaa Shorin algoritmia (\cite{bernstein2017post}). \emph{National Institute of Standards and Technology} eli NIST ylläpitää kvanttiturvallisten julkisen avaimen salausjärjestelmien standardointiprosessia (\cite{alagic2020status}). Kirjoittamishetkellä on käynnissä standardointiprosessin kolmas kierros, johon on valittu seitsemän finalistia. Tutkielmassa käsittelemme kolme kvanttiturvallista salausmenetelmää; koodipohjaisen, hilapohjaisen ja monimuuttujapohjaisen salauksen. Jokaisessa salausmenetelmässä on omat ominaisuutensa ja olivat tasavertaisia keskenään.
 
Tutkielmassa käymme aluksi läpi taustatietoja aiheeseen liittyen. Tutustumme ensin nykyisiin salausjärjestelmiin ja niiden kvanttiturvallisuuteen. Sen jälkeen tutustumme kvanttitietokoneisiin ja kvanttilaskentaan sekä Shorin algoritmiin. Seuraavaksi pääsemme itse aiheeseen, jossa käsittelemme kvanttiturvallista salausta ja perehdymme erilaisiin kvanttiturvallisiin salausmenetelmiin sekä niitä hyödyntäviin julkisen avaimen salausjärjestelmiin. Lopuksi on yhteenveto, jossa käsittelemme tutkimusalueen eli kvanttiturvallisen salauksen nykytilaa ja tulevaisuuden näkymiä.
\chapter{Nykyiset salausjärjestelmät\label{methods}}

Käymme nyt läpi yleisimpiä salausjärjestelmiä, joihin jokainen tietojärjestelmien käyttäjä törmää nykypäivänä tiedostaen tai tiedostamattaan. Käymme läpi salausjärjestelmien ideat ja miten kvanttitietokoneet uhkaavat niitä.

Jatkossa avaimesta puhuttaessa tarkoitetaan tässä kontekstissa salausavainta, jolla voidaan purkaa tai salata lähetettävä viesti. Käytännössä tämä avain on vain \emph{k} bittiä pitkä bittijono, joka on muodostettu jonkin salausavainten luontiin tarkoitetun algoritmin avulla.

\section{Symmetrinen salaus}
 Symmetrisessä salauksessa sekä viestin lähettäjän eli salaajan sekä viestin vastaanottajan eli salauksen purkajan täytyy tietää sama avain (\cite{bellare2005introduction}). Toisin sanoen symmetrisessä salauksessa viestin salaus puretaan samalla avaimella kuin millä viesti salataan.  Symmetrisessä salauksessa ongelmallista on salausavaimen jakaminen turvallisesti ja luottamuksellisesti osapuolien välillä. Jos symmetriseen salaukseen käytetty salausavain joutuu jonkun kolmannen osapuolen käsiin, voi tämä osapuoli purkaa salauksen tällä avaimella. Esimerkki symmetrisestä salausjärjestelmästä on lohkosalaukseen (\emph{blockcipher}) perustuva AES (\emph{Advanced Encryption Standard}) (\cite{117146}). AES jakaa datan 128 bittiä pitkiin lohkoihin ja käyttää 128, 192 ja 256 bittisiä salausavaimia.
 
 Symmetrinen salaus on itsessään jo melko kvanttiturvallinen, koska niiden murtamiseen käytetään yleensä raa'an voiman (\emph{brute-force}) tekniikoita (\cite{mavroeidis2018impact}). Toistaiseksi ei ole löydetty kvanttialgoritmeja, jotka murtaisivat symmetristä salausta käyttäviä järjestelmiä huomattavan tehokkaasti siten, että ne olisi uhattuna muuttua turvattomiksi. Kuitenkin \emph{Groverin algoritmia} hyödyntämällä on mahdollista nopeuttaa kvanttitietokoneella tehtyjä brute-force hyökkäyksiä.
 
 Lyhyesti Groverin algoritmi on kvanttialgoritmi, jota käytetään tiedon hakemiseen (\cite{hayward2008quantum}). Groverin algoritmi löytää \emph{N}-kokoisesta järjestämättömästä listasta halutun alkion $\sqrt{N}$ askeleella. Klassisilla tietokoneilla ja algoritmeilla tällaisessa tehtävässä menee aina vähintään $N$ askelta. Siten Groverin algoritmi tuo neliöllisen nopeutuksen kvanttitietokoneella tehtyihin brute-force hyökkäyksiin (\cite{mavroeidis2018impact}). Kvanttitietokone joka käyttää Groverin algoritmia \emph{n}-bittisen salausavaimen purkamiseen laskee avaimen $\sqrt{2^{n}} = 2^{n/2}$ operaatiolla. Tällöin esimerkiksi 128-bittinen AES-salaus antaisi 64-bittisen suojauksen. Toisaalta näitä hyökkäyksiä voidaan hidastaa esimerkiksi tuplaamalla salausavaimen käyttämä bittien määrä. Lohkosalaukseen perustuvia salausjärjestelmiä AES-192 ja AES-256 pidetään kvanttiturvallisina nyt ja jatkossa (\cite{mavroeidis2018impact}). Nykypäivänä 80-bittistä suojausta pidetään turvallisena rajana klassisella tietokoneella.
 
 \section{Tiivistefunktiot}
 Tiivistefunktiot (\emph{hash functions}) eli puhekielessä hashit  ovat funktioita, jotka ottavat mielivaltaisen pituisen merkkijonon \emph{M} ja muodostavat siitä aina tietyn mittaisen toisen merkkijonon \emph{MD}, jota kutsutaan tiivisteeksi (\cite{sobti2012cryptographic}). Formaalimmin tiivistefunktio on funktio
 \[h: D \to R, \text{ missä } D = \big\{0,1\big\}^{*}, R = \big\{0,1\big\}^{n}, \text{ jollain } n \geq 1.\] Voi siis ajatella, että tällainen hash ottaa merkkijonon ja tiivistää sen toiseksi merkkijonoksi, josta tulee sana tiiviste.
 
 Yleisesti käytettyjä tiivistefunktoita ovat muun muassa SHA1 ja MD-5. Esimerkiksi SHA1 tiivistefunktio muodostaa syötteestä aina 160 bittiä pitkän tiivisteen (\cite{bellare2005introduction}). Tärkeä ominaisuus tiivistefunktioissa on se, että ne toimivat yhteen suuntaan. Viestistä \emph{M} on helposti saatavissa tiiviste \emph{MD} funktiolla \emph{F}, mutta tiivisteestä \emph{MD} ja funktiosta \emph{F} ei voi helposti päätellä tai muodostaa viestiä \emph{M} (\cite{bakhtiari1995cryptographic}). Toisin sanoen kun jostain viestistä, esimerkiksi salasanasta, on tehty tiiviste niin ei ole olemassa yleistä menetelmää, jolla tämä \emph{hashaus} voitaisiin peruuttaa tai saataisiin selville itse funktio \emph{F}.
 
 Tiivistefunktiot ovat symmetrisen salauksen tapaan jo itsessään kvanttiturvallisia (\cite{mavroeidis2018impact}). Tiivistefunktiot muodostavat lohkosalauksen tapaan tietyn mittaisen bittijonon tiivisteen muodossa. Groverin algoritmia hyödyntäen on mahdollista tehdä hyökkäys tiivistefunktioita vastaan löytämällä törmäys (\emph{collision}) kahden tiivisteen välillä. Törmäyksellä tarkoitetaan tilannetta milloin kaksi eri syötettä antavat saman tiivisteen. Groverin algoritmilta voi suojautua pidentämällä tiivisteen pituutta. Turvallinen raja on, että \emph{n}-bitin suojaukseen tarvitaan 3\emph{n}-bittiä pitkä tiiviste. Esimerkiksi SHA-2 ja SHA-3 standardin tiivistefunktiot pitkillä bittijonoilla ovat kvanttiturvallisia (\cite{mavroeidis2018impact}).
 
 \section{Julkisen avaimen salaus}
 Epäsymmetrisessa eli julkisen avaimen salauksessa on käytössä kaksi avainta. Toinen on niin kutsuttu julkinen avain (\emph{public key}) ja toinen salainen avain (\emph{secret key\footnote{Joskus myös \emph{private key}}}) (\cite{bellare2005introduction}). Julkisen avaimen salausjärjestelmässä jokainen käyttäjä luo itselleen parin avaimia, jotka ovat salausavain ja purkuavain. Tässä tapauksessa purkuavain on salainen avain ja sitä ei missään nimessä pidä paljastaa muille. Tällöin salausavain on julkinen avain, joka saa ja pitää olla julkisesti esillä. 
 
 Olkoon meillä lähettäjä \emph{S}, joka haluaa lähettää salatun viestin vastaanottajalle \emph{R}. Viestin lähetys toimii nyt seuraavasti:
 \begin{enumerate}
     \item \emph{S} salaa viestin \emph{R}:n julkisella avaimella $PK_{R}$.
     \item \emph{S} lähettää viestin vastaanottajalle \emph{R}.
     \item \emph{R} saa \emph{S}:n viestin vastaan.
     \item \emph{R} purkaa \emph{S}:n salauksen omalla salaisella avaimellaan $SK_{R}$.
 \end{enumerate}
Ero symmetriseen salaukseen epäsymmetrisessa salauksessa on siis se, että salaus puretaan eri avaimella kuin millä se salataan.
 
  Julkisen avaimen salausjärjestelmien tarjoama suojaus perustuu tiettyihin vaikeisiin lukuteorian ongelmiin (\cite{mavroeidis2018impact}). Monissa ongelmissa kyse on matemaattisista funktioista, jotka on helppo laskea yhteen suuntaan, mutta niiden peruuttaminen tai käänteisoperaatio taas on hyvin hankalaa. Näitä ongelmia ovat muun muassa suurten alkulukujen tekijöihinjako (\emph{Integer Factorization Problem}) sekä niin kutsuttu diskreetin logaritmin ongelma (\emph{Discrete Logarithm Problem, DLP}).
  Yksi ensimmäisiä ja tänäkin päivänä käytetyimpiä julkisen avaimen salausjärjestelmiä on RSA, joka on nimetty keksijöidensä (Rivest, Shamir ja Adleman) mukaan (\cite{montgomery1994survey}). RSA:n tarjoama suojaus perustuu suurten alkulukujen tekijöihin jakoon. RSA:n avainpari luodaan seuraavasti  (\cite{10.1145/359340.359342}):
  
  \begin{enumerate}
  
      \item Valitaan kaksi suurta satunnaista alkulukua \emph{p} ja \emph{q}, siten että $p \neq q$. Nämä luvut $p$ ja $q$ on pidettävä salassa.
      
      \item Lasketaan $n = pq$. Tästä saatava $n$ saa olla julkisesti esillä.
      
      \item Valitaan jokin suuri satunnainen kokonaisluku $d$ siten että,
      \[syt(d, (p-1)(q-1)) = 1\]
      Tarkemmin sanottuna $d$ ja $(p-1)(q-1)$ ovat nyt keskenään suhteellisia alkulukuja (\emph{coprime} tai \emph{relatively prime}). Tämä tarkoittaa sitä, että niiden suurin yhteinen tekijä eli suurin kokonaisluku, joka jakaa sekä $d$:n että $(p-1)(q-1)$:n on 1.
      
      \item Lasketaan $e$, joka on $d$:n käänteisluku
      \[e d \equiv 1 \pmod{(p-1) (q-1)}\]
  \end{enumerate}
  Salausavain on nyt pari positiivisia kokonaislukuja ($e$, $n$). Purkuavain on pari positiivisia kokonaislukuja ($d$, $n$).
  
 Esitellään seuraavaksi diskreetin logaritmin ongelma. Olkoon meillä jokin äärellinen kunta $GF$($q$), jossa on $q$ alkiota (\cite{1055638}). Olkoon
 \begin{equation}
     \label{eq10}
     Y = \alpha^{X} \mod q, \; \text{jollekkin} \; 1 \leq X \leq q-1
 \end{equation}
 missä $\alpha$ on jokin määrätty alkio $GF$($q$):sta. Yhtälö voidaan muuttaa muotoon siten, että ratkaistaan $X$ 
 \begin{equation}
     \label{eq11}
     X = \log_{\alpha} Y \mod q, \; \text{jollekkin} \; 1 \leq Y \leq q-1
 \end{equation}
 Nyt siis $Y$:n selvittäminen $X$:stä on helppoa ja vaatii enimmillään $2 \times \log_{2}q$ kertolaskuoperaatiota. Toisaalta $X$:n selvittäminen $Y$:stä ei ole helppoa ja vaatii parhailla algoritmeilla $q^{1/2}$ kertolaskuoperaatiota. Esimerkiksi julkisen avaimen salausjärjestelmät kuten Diffie-Hellman ja elliptisen käyrän salaus (\emph{Elliptic Curve Cryptography, ECC}) hyödyntävät tätä (\cite{mavroeidis2018impact}).
 
 Edellä mainitut suurien kokonaislukujen tekijöihinjako ja diskreetin logaritmin löytäminen ovat ongelmia, joita on tutkittu paljon, mutta siitä huolimatta niille ei ole löytynyt polynomisessa ajassa toimivia algoritmeja (\cite{doi:10.1137/S0036144598347011}). Peter Shor esitteli vuonna 1994 Shorin algoritmin (\cite{mavroeidis2018impact}), jota käyttäen edellä mainitut vaikeat lukuteorian ongelmat voidaan ratkaista tehokkaasti eli polynomisessa ajassa kvanttitietokoneella (\cite{doi:10.1137/S0036144598347011}). Käsittelemme Shorin algoritmia myöhemmin kappaleessa 3.2 kvanttilaskennan yhteydessä.
 
 \section{Digitaalinen allekirjoitus}
 Digitaalinen allekirjoitus (\emph{digital signature}) on menetelmä, jolla voidaan todentaa digitaalisia dokumentteja (\cite{1649003}). Sen toimintaperiaate on analoginen perinteisiin allekirjoituksiin verrattuna. Digitaalisella allekirjoituksella todennetaan muun muassa digitaalisen dokumentin datan eheys. Eheydellä voidaan tässä kontekstissa tarkoittaa esimerkiksi dokumentin päivämäärää tai sisältämää informaatiota.
 
 Digitaalinen allekirjoitus on käytännössä bittijono, jonka jokin digitaalisen allekirjoituksen luomiseen tarkoitettu algoritmi luo. Tämä allekirjoitus muodostetaan allekirjoitettavien dokumenttien ja jonkin vain viestin lähettäjän määrittelemän tiedon mukaan. Digitaalisen allekirjoituksen luominen koostuu yleensä kolmesta eri vaiheesta; salausavaimen generointi, digitaalisen allekirjoituksen generointi ja digitaalisen allekirjoituksen varmennus.
 
 Edellä mainittua julkisen avaimen salausta voidaan käyttää digitaalisen allekirjoituksen luomiseen (\cite{mavroeidis2018impact}). Esimerkiksi dokumentin allekirjoittaja \emph{A} voi allekirjoittaa dokumentin omalla salaisella avaimellaan $SK_{A}$. Dokumentin varmentaja \emph{V} voi siten todentaa \emph{A}:n allekirjoituksen \emph{A}:n julkisella avaimella $PK_{A}$. Kaksi suosituinta ja eniten käytettyä julkisen avaimen salaukseen perustuvaa digitaalisen allekirjoittamisen menetelmää ovat RSA\footnote{Tässä RSA:lla tarkoitetaan samannimistä digitaalisen allekirjoituksen menetelmää, joka hyödyntää kappaleessa 2.3 mainittua RSA-algoritmia.} ja DSA (\emph{Digital Signature Algorithm}) (\cite{1649003}). Koska RSA:n turvallisuus perustuu suurten alkulukujen tekijöihinjakoon ja DSA:n turvallisuus diskreettien logaritmien laskemiseen, niin Shorin algoritmi uhkaa myös näitä digitaalisen allekirjoituksen menetelmiä.
\chapter{Kvanttilaskenta}

\section{Kvanttilaskenta ja kvanttitietokoneet}
Klassisessa eli niin sanotussa tavallisessa tietokoneessa informaatio esitetään bitteinä ja yksi bitti on tiedon perusyksikkö (\cite{doi:10.1080/23742917.2016.1226650}). Yksittäinen bitti voi saada vain yhden diskreetin arvon kerrallaan ja näitä arvoja kuvataan yleensä arvoilla 0 ja 1. Bitti voi siis olla arvoltaan vuorostaan joko 0 tai 1. Klassisessa tietokoneessa näitä bittejä käsitellään loogisten porttien kautta.

Kvanttitietokoneen laskennan perusyksikkö taas on kubitti eli kvanttibitti. Kvanttitietokoneessa kubitteja käsittelee kvanttiportit kuten klassisessa tietokoneessa bittejä käsittelee loogiset portit. Kubitti voi olla bitin tavoin myös kahdessa eri tilassa. Merkitään näitä tiloja symboleilla
$\ket{0}$ ja $\ket{1}$.\footnote{Kvanttimekaniikassa notaatiota $\ket{}$ kutsutaan Diracin notaatioksi. Se on yleinen notaatio kvanttimekaniikassa tarkoittaen tilaa (\cite{nielsen2001quantum}).} Kubitti eroaa tavallisesta bitistä siten, että se voi olla tilojen $\ket{0}$ tai $\ket{1}$ sijaan niin sanotussa superpositiotilassa. Tämä tarkoittaa siis sitä, että mittauksen aikana kubitti voi olla arvoltaan 0 tai 1, mutta laskennan aikana se on yhtäaikaa sekä 0 että 1. Eli tilat $\ket{0}$ ja $\ket{1}$ voivat esiintyä yhtäaikaa. Superpositiotilassa olevalle kubitille tehty operaatio vaikuttaa kumpaankin arvoon samanaikaisesti (\cite{mavroeidis2018impact}). Kubitit voivat myös lomittua keskenään, jolloin puhutaan kvanttilomittumisesta (\emph{Quantum entanglement}). Kun kaksi kubittia lomittuu keskenään, niin sen jälkeen niiden tiloja ei voi enää kuvailla riippumatta toisistaan. Sen sijaan ne kuvataan yhtenä objektina, jolla on neljä erillistä tilaa.

Kvanttitietokoneen tehokkaamman laskentatehon mahdollistaa siis nämä edellä mainitut kvanttimekaniikan ilmiöt, kvanttisuperpositio ja kvanttilomittuminen. Kvanttitietokoneella on siis mahdollista suorittaa laskentaa samanaikaisesti eli rinnakkain jokaiselle kubitin eri tilalle (\cite{doi:10.1080/23742917.2016.1226650}). Tätä kutsutaan kvanttirinnakkaisuudeksi (\emph{Quantum parallelism}). Kvanttitietokone joka siis käyttää laskentaan $n$ kubittia mahdollistaa samanaikaisesti $2^{n}$ rinnakkaisesti suoritettavaa operaatiota (\cite{mavroeidis2018impact}). Esimerkiksi kolmella kubitilla laskentaa tekevä kvanttitietokone mahdollistaa 8 eri tilaa samanaikaisesti. Kolmella bitillä laskentaa tekevä klassinen tietokone mahdollistaa myös 8 eri tilaa, mutta vain yksi näistä kahdeksasta tilasta voi olla kerrallaan voimassa.  Kvanttitietokone tarjoaa siis eksponentiaalisesti nopeamman laskentatehon. Tätä eksponentiaalista laskentatehoa hyödyntävät algoritmit uhkaavat muun muassa julkisen avaimen salausta.

Täytyy kuitenkin huomata, että kvanttilaskennan tuoma eksponentiaalinen laskentateho ei ole itsestäänselvyys. Kvanttitietokoneet ovat ylivoimaisia klassisiin tietokoneisiin verrattuna vasta silloin kun ongelmaan voidaan soveltaa algoritmia, joka osaa hyödyntää kvanttitietokoneiden mahdollistamaa kvanttirinnakkaisuutta. Esimerkiksi kertolaskujen laskemisessa kvanttitietokone ei juuri tuo huomattavaa etua verrattuna klassisiin tietokoneisiin.

Nykypäivän kvanttitietokoneet ovat vielä kehitysvaiheessa (\cite{doi:10.1080/23742917.2016.1226650}). Ne vaativat suuren ja kalliin infrastruktuurin ympärilleen muun muassa jäähdytystä ja sähkömagneettiselta säteilyltä suojautumista varten. Kuitenkin tilanne oli aikoinaan samanlainen kun klassiset tietokoneet alkoivat ilmaantumaan. Tulevaisuudessa on siis odotettavissa, että kvanttitietokoneiden fyysinen koko ja vaatima infrastruktuuri tulee pienentymään.

\section{Shorin algoritmi}
Kappaleessa 2 nykyisistä salausjärjestelmistä puhuttaessa mainitsimme Shorin algoritmin. Vuonna 1994 Peter Shor esitteli Shorin algoritmin, jolla voidaan jakaa kokonaisluku tekijöihinsä kvanttitietokoneella tehokkaasti eli polynomisessa ajassa (\cite{hayward2008quantum}). Shorin algoritmin toiminta ja tehokkuus perustuu siihen kuinka se hyödyntää edellä mainittua kvanttirinnakkaisuutta. Tämä tarkoittaa siis sitä, että kvanttitietokoneella Shorin algoritmi suorittaa eksponentiaalisen määrän operaatioita yhdessä askeleessa. Klassisella tietokoneella sama laskenta veisi eksponentiaalisen määrän aikaa.

Shor huomasi, että Eukleideen kehittämää tekijöihinjakoalgoritmia pystyttiin nopeuttamaan käyttäen kvanttirinnakkaisuutta apuna (\cite{doi:10.1080/23742917.2016.1226650}). Kuvaillaan seuraavaksi Eukleideen algoritmin toimintaa (\cite{hayward2008quantum}). Olkoon meillä funktio
\begin{equation}
    \label{eq1}
    F(a) = x^{a} \mod n
\end{equation}

joka on jaksollinen kun $x$ on suhteellinen alkuluku\footnote{ks. kappale 2.2.} $n$:n kanssa ja $n$ on nyt luku jonka tahdomme jakaa tekijöihin. Koska $F(a)$ on jaksollinen funktio\footnote{Jaksollinen funktio tarkoittaa siis funktiota, joka toistuu samanlaisena tietyn jakson välein.}, niin sillä on tällöin jokin jakso $r$. Shorin algoritmi siis yrittää löytää $r$:n eli funktion $F(a) = x^{a} \mod n$ jakson. Tiedetään, että 
\begin{equation}
    \label{eq2}
    \begin{split}
    & x^{0} \mod n = 1 \; \text{koska} \;  x^0 = 1 \; \text{joten} \\
    & x^{r} \mod n = 1 \; \text{ja} \\
    & x^{2r} \mod n = 1 \\
    & \; \text{ja niin edelleen kun $r$ saa monikertoja.}
    \end{split}
\end{equation}
Siten voidaan johtaa
\begin{equation}
    \label{eq3}
    \begin{split}
    x^{r} & \equiv 1 \mod n \\
    (x^{r/2})^{2} = x^{r} & \equiv 1 \mod n \\
    (x^{r/2})^{2} - 1 & \equiv 0 \mod n
    \end{split}
\end{equation}
ja jos $r$ on parillinen niin
\begin{equation}
    \label{eq4}
    (x^{r/2} - 1)(x^{r/2} + 1) \equiv 0 \mod n
\end{equation}

Nyt $(x^{r/2} - 1)(x^{r/2} + 1)$ on $n$:n eli luvun jonka haluamme jakaa tekijöihin monikerta. Joten jos $\lvert x^{r/2} \rvert$ ei ole 1, niin ainakin toisella $(x^{r/2} - 1), (x^{r/2} + 1)$ on sama tekijä kuin $n$:llä. Nyt laskemalla $syt(x^{r/2} - 1, n)$ ja $syt(x^{r/2} + 1, n)$ saadaan $n$:n tekijä\footnote{ks. kappale 2.2}.

Simuloidaan nyt Eukleideen algoritmia pienellä syötteellä (\cite{doi:10.1080/23742917.2016.1226650}). Olkoon $n$ = 15 ja haluamme jakaa sen tekijöihin.

\begin{enumerate}
    \item Valitaan suhteellinen alkuluku. Valitaan 11, koska $syt(11,15) = 1$.
    
    \item Ryhdytään laskemaan $x^{a} \equiv 1 \mod n$, missä $a \in \mathbb{N}$.
    
    \begin{enumerate}
    
        \item $11^{1} \equiv 1 \mod 15 = 0$ kokonaista ja jakojäännöksenä 11.
        
        \item $11^{2} \equiv 1 \mod 15 = 121$ eli 8 kokonaista ja jakojäännöksenä 1.
        
        \item $11^{3} \equiv 1 \mod 15 = 1331$ eli 88 kokonaista ja jakojäännöksenä 11.
    \end{enumerate}
    
    \item Jatketaan näin korottamalla $x = 11$ aina korkeampiin potensseihin. Jakojäännös jaettaessa luvulla 15 vaihtelee antaen joko luvun 1 tai 11.
    
    \item Saamme nyt jaksoksi $r = 2$, koska kun lukua 11 jaetaan luvulla 15 saamme aina jakojäännökseksi joko 1 tai 11.
    
    \item Lasketaan nyt $x^{r/2}$, jolloin saadaan $x^{r/2} = 11^{2/2} = 11$.
    
    \item Lasketaan nyt $(x^{r/2}-1)$ ja $(x^{r/2}+1)$, jolloin saadaan $(11^{2/2} - 1) = (11 - 1) = 10$ ja $(11^{2/2} + 1) = (11 + 1) = 12$.
    
    \item Lasketaan $syt(10,15)$ ja $syt(12,15)$, jolloin saadaan $syt(10,15) = 5$ ja $syt(12,15) = 3$ ja huomataan, että 3 ja 5 ovat alkuluvut jotka muodostavat luvun 15.
    
\end{enumerate}

Ongelmana nyt on, että kun $n$ kasvaa hyvin suureksi, niin jakson $r$ selvittäminen funktiolla $F(a)$ käy hyvin hankalaksi. Klassisilla tietokoneilla tämän ratkaisemiseen käytetään apuna nopeita Fourier-muunnoksia (\emph{Fast-Fourier Transform, FFT}), mutta suurilla $n$:n arvoilla niiden laskeminen on hyvin epäkäytännöllistä. Shorin algoritmi ratkaisee jakson $r$ etsimisen tehokkaasti käyttäen apunaan kvanttirinnakkaisuutta. Lisäksi jakson $r$ määrittäminen on ainut kohta algoritmissa, jossa tarvitaan kvanttitietokonetta ja missä kvanttilaskenta tuo huomattavaa etua.

Nopein suurten kokonaislukujen tekijöihinjako algoritmi on yleinen lukukuntaseula (\emph{General Number Field Sieve, GNFS}) (\cite{doi:10.1137/S0036144598347011}). \emph{GNFS}:n asymptoottinen ajoaika on $exp(c(\log n)^{1/3} (\log \log n)^{2/3})$, missä $n$ on kokonaisluku, joka halutaan jakaa tekijöihin ja $c$ on vakio. GNFS on siis ajan suhteen eksponentiaalinen algoritmi. Shorin algoritmin asymptoottinen ajoaika kvanttitietokoneella on $O((\log n)^{2}(\log \log n) (\log \log \log n))$, missä $n$ on luku, joka halutaan jakaa tekijöihin. Tämän lisäksi tarvitaan kvanttitietokoneella tehdyn laskennan jälkeen $O(\log n)$ verran prosessointia klassisella tietokoneella. Tämä tehdään, jotta kvanttitietokoneen antama tulos eli $n$:n tekijät voidaan muuttaa klassisen tietokoneen ymmärtämään muotoon.

Tähän saakka Shorin algoritmi on saatu skaalautumaan seitsemällä kubitilla (\cite{doi:10.1126/science.aad9480}). 1024-bittisen RSA-avaimen tekijöiden jakoon tarvittaisiin 2000-kubittinen kvanttitietokone (\cite{mavroeidis2018impact}).
\include{Ch.40_Kvanttiturvallinensalaus}
\chapter{Yhteenveto\label{conclusions}}
% Puhu mitä työssä on tehty!

Tutkielmassa tutustuimme pintapuolisesti nykyisiin salausjärjestelmiin ja niiden kvanttiturvallisuuten, kvanttilaskentaan ja Shorin algoritmiin sekä kvanttiturvallisiin salausjärjestelmiin, joihin Shorin algoritmia ei voi soveltaa ja jotka perustuvat kvanttitietokoneellekkin vaikeisiin ongelmiin.

Kvanttitietokoneiden kehittyessä kvanttiturvallinen salaus on tullut ajankohtaiseksi. Kvanttitietokoneiden mahdollistama eksponentiaalinen laskentateho klassisiin tietokoneisiin verrattuna uhkaa nykyisiä julkisen avaimen salausjärjestelmiä. Etenkin kvanttirinnakkaisuutta hyödyntävä Shorin algoritmi ratkaisee klassiselle tietokoneelle vaikeat laskennalliset ongelmat tehokkaasti. Esimerkiksi yksi ensimmäisiä ja käytetyimpiä julkisen avaimen salausjärjestelmiä, RSA, käyttää suurten kokonaislukujen tekijöihinjakoa salauksessaan minkä Shorin algoritmi pystyy murtamaan tehokkaasti.

Tutkielmassa keskityimme kolmeen eri kvanttiturvalliseen salausmenetelmään; koodipohjaiseen, hilapohjaiseen ja monimuuttujapohjaiseen salaukseen. Jokainen salaus perustui omaan NP-täydelliseen ongelmaan; koodipohjaisella  \emph{Generic Decoding Problem}, hilapohjaisella \emph{Shortest Vector Problem} sekä muut hiloihin liittyvät laskennalliset ongelmat ja monimuuttujapohjaisella \emph{MQ}-ongelma. Jokaisessa salauksessa oli omat ominaisuutensa ja kehityskohteensa. Jokainen salausmenetelmä todettiin laskennallisesti tehokkaaksi ja nopeaksi, mutta monimuuttujapohjainen salaus todettiin kaikista kolmesta nopeimmaksi. Koodipohjaisesta ja monimuuttujapohjaisesta salauksesta todettiin huonona puolena suuret julkiset avaimet, jotka rajoittavat näiden salausjärjestelmien käyttöä esimerkiksi sulautetuissa järjestelmissä. Lisäksi koodipohjaisessa ja monimuuttujapohjaisessa oli puutteita kvanttiturvallisuuden teoriapohjassa, vaikka käytännössä ovat pysyneet kvanttiturvallisina. Hilapohjaisesta salauksesta todettiin kehityskohteena salatun viestin koon kasvaminen monikymmenkertaiseksi. Esittelimme myös jokaista salausmenetelmää edustavan NIST:in standardointikierroksen finalistin. Finaalissa ovat tällä hetkellä julkisen avaimen salausmenetelmiä edustamassa koodipohjainen \emph{Classic McEliece} ja hilapohjainen \emph{NTRU}. Allekirjoittamiseen luotua järjestelmää edustaa monimuuttujapohjainen \emph{Rainbow}.

Vaikka suuren kokoluokan kvanttitietokoneet eivät juuri tällä hetkellä ole suuri uhka tietoturvalle niin kvanttiturvalliseen salaukseen pitäisi varautua jo nyt. Retrospektiivisellä purkamisella tarkoitetaan hyökkäystä, missä hyökkääjä kerää salattua tietoa talteen tulevaisuutta varten ja purkaa salatun tiedon jälkeenpäin kun löydetään keino purkaa salaus.

%%%%%%%%%%%%%%%%%%%%%%%%%%%%%%%%%%%%%%%%%%%%%%%%%%%%%%%%%
\cleardoublepage                          %fixes the position of bibliography in bookmarks
\phantomsection
\addcontentsline{toc}{chapter}{\bibname}  % This lines adds the bibliography to the ToC
\printbibliography

%%%%%%%%%%%%%%%%%%%%%%%%%%%%%%%%%%%%%%%%%%%%%%%%%%%%%%%%%
\backmatter
\begin{appendices}

%% A sample Appendix
%
\appendix{Sample Appendix\label{appendix:sample}}

You can add one or more appendices to your thesis.

%% another appendix
%
\appendix{Instructions for LaTex}

\section{General Setup}

In the HY-CS-main.tex file you will find the following STEPS 0--5. Below you can find related instructions.
\vspace{0.5cm}

\textbf{STEP 0 -- Access the thesis template}

\begin{itemize}
\item Import the thesis template into a new Overleaf project. The easiest way to do it is to:
\begin{itemize}
    \item Obtain a zip file of the LaTeX template from the webpage of your programme.
    \item Go to \url{https://www.overleaf.com/edu/helsinki} and login to Overleaf with your university credentials.
    \item Go to the list of your projects at \url{https://www.overleaf.com/project}, click ``New Project'' and ``Upload Project''.,  the projects under your account 
    \item Then upload the zip with the template.
    \item You are now ready to write your thesis in Overleaf by editing the template, you can start by renaming the project.
\end{itemize}
\end{itemize}


{\textbf{STEP 1 -- BSc or MSc thesis?}}
\begin{enumerate}
\item Select whether your are writing BSc (tkt) or MSc (csm for CS) thesis.
\item Select your language: \texttt{finnish}, \texttt{english}, or \texttt{swedish}.
\item If you are writing MSc select your line / track.
\end{enumerate}


{\textbf{STEP 2 -- Set up your personal information}}

\begin{enumerate}
\item Specify the title of your thesis with \texttt{\textbackslash title\{\}}.
\item Specify your name to the author field with \texttt{\textbackslash author\{\}}.
\item Specify the names of your supervisors of the thesis with \texttt{\textbackslash supervisors\{\}}.
\item Specify the keywords of the thesis with \texttt{\textbackslash keywords\{\}}.
\item Specify the ACM classification terms of the thesis with \texttt{\textbackslash classification\{\}}. See \url{https://dl.acm.org/ccs} for more information.
\end{enumerate}

{\textbf{STEP 3 -- Write your abstract}}

\begin{itemize}
\item You can have the abstract in multiple languages with the \texttt{otherlanguages} environment. The example below shows how to provide an English abstract: 

\begin{verbatim}
\begin{otherlanguage}{english} 
\begin{abstract}
Your abstract text goes here. 
\end{abstract} 
\end{otherlanguage}
\end{verbatim}

\end{itemize}

{\textbf{STEP 4 -- Writing your thesis}}

\begin{enumerate}
\item There are some minimal contents and instructions below 
\item Remove, or comment out, this appendix from your thesis.
\end{enumerate}

{\textbf{STEP 5 -- Set your bibliography style}}

\begin{itemize}
\item The default is Author-Year style (Einstein, 1905), but it can be easily changed to numbered [1] or alphabetical [Ein05] , as the examples of these are in comments.
\item Discuss the style to use with your supervisor.
\end{itemize}

\section{Bibliography in Latex}

The bibliography is defined in a separate \texttt{.bib} file. For this template, it is named \texttt{bibliography.bib} and includes the content show in Figure~\ref{bibexamples}.

Chapter Bibliography lists all the works that you refer to in your text. You refer to the works in the bibliography using an appropriate \emph{citation key}.
%
%This thesis template contains an example of a bibliography.


References are done using \texttt{\textbackslash citep\{einstein\}}, which generates in text a citation formatted according to the selected style \citep{einstein}, or \texttt{\textbackslash citep\{latexcompanion,knuth99\}}, which generates \citep{latexcompanion,knuth99}. 
As examples of a different kinds of citations (see how these look in the Latex source), we can write \citep{einstein} to refer to the work written by \citeauthor{einstein} in \citeyear{einstein}, because the work by \citet{einstein} appears in the bilbliography included in this template.

Note that there are different possible styles for the bibliography and citation keys.
%
Consult your supervisors on the chosen style -- and once you arrive at a preferred style, use it consistently throughout the thesis.

\begin{figure}[ht]
    \centering
    \begin{scriptsize}
\begin{verbatim}
@article{einstein,
    author =       "Albert Einstein",
    title =        "{Zur Elektrodynamik bewegter K{\"o}rper}. ({German})
        [{On} the electrodynamics of moving bodies]",
    journal =      "Annalen der Physik",
    volume =       "322",
    number =       "10",
    pages =        "891--921",
    year =         "1905",
    DOI =          "http://dx.doi.org/10.1002/andp.19053221004"
}
 
@book{latexcompanion,
    author    = "Michel Goossens and Frank Mittelbach and Alexander Samarin",
    title     = "The \LaTeX\ Companion",
    year      = "1993",
    publisher = "Addison-Wesley",
    address   = "Reading, Massachusetts"
}

@book{knuth99,
    author    = "Donald E. Knuth",
    title     = "Digital Typography",
    year      = "1999",
    publisher = "The Center for the Study of Language and Information",
    series    = "CLSI Lecture Notes (78)"
}\end{verbatim}
\end{scriptsize}
    \caption{Examples of bibliographic reference in .bib file.}
    \label{bibexamples}
\end{figure}

%In the last reference url field the code \verb+%7E+ will translate into \verb+~+ once clicked in the final pdf.

\section{Some instructions about writing in Latex}

The following gives some superficial instructions for using this template for a Master's thesis. For guidelines on thesis writing you can consult various sources, such as university courses on scientific writing or your supervisors.

For more detailed instructions, just google, e.g., "Overleaf table positioning", and your chances of finding good info are pretty good.  


\section{Figures}
Besides text, here are simple examples how you can add figures and tables in your thesis.
Remember always to refer to each figure in the main text and provide them with a descriptive caption.

Figure~\ref{fig:logo} is an example of a figure in the document (see the source about how to add them). 
%Using figures is particularly useful to display plots of experimental results.

\begin{figure}[ht] 
\begin{center}
\includegraphics[width=0.3\textwidth]{template/figures/HY-logo-ml.png}
\caption{University of Helsinki flame-logo for Faculty of Science.\label{fig:logo}}
\end{center}
\end{figure}

\section{Tables}

Table~\ref{table:results} gives an example of a table.
Remember always to cite the table in the main text, table captions go on top of the table. 

\begin{table}[h] % h positions the table here, t! would force on top of the page, or example.
\begin{center}
\caption{Experimental results.\label{table:results}} % caption is here to make it on top
\begin{tabular}{l||l c r} 
Experiment & 1 & 2 & 3 \\ 
\hline \hline 
$A$ & 2.5 & 4.7 & -11 \\
$B$ & 8.0 & -3.7 & 12.6 \\
$A+B$ & 10.5 & 1.0 & 1.6 \\
\hline
%
\end{tabular}
\end{center}
\end{table}



%% yet another appendix
%
\appendix{Tutkielmapohjan käyttöohjeet}
\label{appendix:instructions_finnish}

\section{Ensiaslkeleet}

\texttt{HY-CS-main.tex} tiedosto sisältää viisi askelta STEPS 0--5. Alla on kuvattu, mitä nämä askeleet tarkoittavat ja miten niitä seuraamalla luot pohjan tutkielmallesi.
\vspace{0.5cm}

\textbf{STEP 0 -- Kopioi tutkielmapohja}

\begin{itemize}
\item Hae tutkielmapohja uuteen Overleaf-projektiin. Tämä käy helpoiten seuraavasti:
\begin{itemize}
    \item Lataa Latex-pohjan zip-tiedosto koulutusohjelman sivuilta.
    \item Mene osoitteeseen \url{www.overleaf.com/edu/helsinki} ja kirjaudu Overleafiin yli\-opiston tunnuksillasi.
    \item Overleafissa (\url{https://www.overleaf.com/project}), klikkaa ``New Project'' and ``Upload Project''.
    \item Valitse lataamasi tutkielmapohjan zip-tiedosto.
    \item Nyt voit lähteä kirjoittamaan tutkielmaasi suoraan pohjaan, voit aloittaa esim. vaihtamalla projektin nimen.
\end{itemize}
\end{itemize}


{\textbf{STEP 1 -- BSc vai MSc tutkielma?}}
\begin{enumerate}
\item Valitse (tiedostossa \texttt{HY-CS-main.tex}) oletko tekemässä BSc (tkt) vai MSc (csm tietojenkäsittely) tutkielmaa.
\item Valitse kieli jolla kirjoitat tutkielman: \texttt{finnish}, \texttt{english} tai \texttt{swedish}.
\item Jos olet kirjoittamassa maisterintutkielmaa, valitse linja/opintosuunta.
\end{enumerate}


{\textbf{STEP 2 -- Aseta henkilökohtaiset tietosi}}

\begin{enumerate}
\item Kirjoita alustava otsikko tutkielmallesi: \texttt{\textbackslash title\{\}}.
\item Kirjoita oma nimesi kohtaan \texttt{\textbackslash author\{\}}.
\item Lisää ohjaajien nimet \texttt{\textbackslash supervisors\{\}}.
\item Määrittele avainsanat \texttt{\textbackslash keywords\{\}}.
\item Määritä tutkielmasi ACM luokittelutermit \texttt{\textbackslash classification\{\}}. Ks. lisätietoa: \url{https://dl.acm.org/ccs}.
\end{enumerate}

{\textbf{STEP 3 -- Kirjoita tiivistelmä}}

%\begin{itemize}
%\item 
Voit kirjoittaa tiivistelmän (koko tiivistelmäsivu) eri kielillä \texttt{otherlanguages}-ym\-pä\-ris\-tön avulla. Alla esimerkki jolla kirjoitat englanninkielisen tiivistelmän muulla kuin englannin kielellä kirjoitettuun tutkielmaan:

\begin{verbatim}
\begin{otherlanguage}{english} 
\begin{abstract}
Your abstract text goes here. 
\end{abstract} 
\end{otherlanguage}
\end{verbatim}

%\end{itemize}

{\textbf{STEP 4 -- Kirjoita tutkielma}}

\begin{enumerate}
\item Kirjoittamisesta Latexilla löydät hieman ohjeita alempaa.
\item Poista tämä liite ja muu ohjeistus tutkielmastasi, esim. kommentoimalla.
\end{enumerate}

{\textbf{STEP 5 -- Aseta kirjallisuuslähdeluettelon tyyli}}

\begin{itemize}
\item Oletustyylin tekijä-vuosi, eli (Einstein, 1905), voit vaihtaa viittaustyylin (tiedostossa \texttt{HY-CS-main.tex}) helposti (eri mallit kommentoituna) esim. numeroituun [1], tai aakkostyyliin [Ein05].
Lisää ohjeita liittyen viittaustyylin säätämiseen {Bib}\TeX issä löytyy verkosta: \url{https://ctan.org/pkg/biblatex}
\item Sovi käytettävä tyyli ohjaajasi kanssa. 
\end{itemize}

\section{Kirjallisuusviitteet Latexissa}

Kirjallisuuslähteet ylläpidetään erillisessä .bib-tiedostossa. Tässä tutkielmapohjassa käy\-te\-tyt kirjallisuuslähteet, joista esimerkkejä kuvassa~\ref{bibexamples-fi}, löytyvät tiedostosta\newline \texttt{bibliography.bib}.

\begin{figure}[ht]
    \centering
    \begin{scriptsize}
\begin{verbatim}

@article{einstein,
    author =       "Albert Einstein",
    title =        "{Zur Elektrodynamik bewegter K{\"o}rper}. ({German})
        [{On} the electrodynamics of moving bodies]",
    journal =      "Annalen der Physik",
    volume =       "322",
    number =       "10",
    pages =        "891--921",
    year =         "1905",
    DOI =          "http://dx.doi.org/10.1002/andp.19053221004"
}

@book{latexcompanion,
    author    = "Michel Goossens and Frank Mittelbach and Alexander Samarin",
    title     = "The \LaTeX\ Companion",
    year      = "1993",
    publisher = "Addison-Wesley",
    address   = "Reading, Massachusetts"
}

@book{knuth99,
    author    = "Donald E. Knuth",
    title     = "Digital Typography",
    year      = "1999",
    publisher = "The Center for the Study of Language and Information",
    series    = "CLSI Lecture Notes (78)"
}
\end{verbatim}
\end{scriptsize}
    \caption{Esimerkkejä kirjallisuuslähteiden kuvaamisesta .bib-tiedostossa.}
    \label{bibexamples-fi}
\end{figure}

Viitteet kirjallisuuslähteisiin muodostetaan komennolla \texttt{\textbackslash citep\{einstein\}}, josta generoituu tekstiin valitun viittaustyylin mukaisesti muotoiltu viite \citep{einstein}, tai \texttt{\textbackslash citep\{latexcompanion,knuth99\}}, josta tekstiin puolestaan generoituu \citep{latexcompanion,knuth99}. 
Voit esimerkiksi kirjoittaa \citep{einstein} viitataksesi julkaisuun, jonka on kirjoittanut \citeauthor{einstein} vuonna \citeyear{einstein}, kun vain lähde \citet{einstein} on oikein lisättynä kirjallisuuslähdetiedostossa (katso miltä nämä näyttävät Latex lähdekoodissa).

Tekstissä viitatut kirjallisuuslähteet tulevat automaattisesti viiteluetteloon. Kirjallisuuslähteiden tietojen oikeellisuus ja yhdenmukaisuus .bib-tiedostossa vaikuttavat luonnollisesti siihen, miten tiedot tutkielmassa näyttäytyvät. Tämä on syytä huomioida, sillä esim.\ verkosta valmiiksi {Bib\TeX} muodossa löytyvien tietojen täydellisyyten tai samanmuotoisuuteen ei pidä sokeasti luottaa.  


Keskustele viittaustyylin valinnasta ohjaajan kanssa. 
%Joitain vaihtoehtoja on osoitteessa:\\ 
%\url{https://www.overleaf.com/learn/latex/Biblatex_bibliography_styles}.
%\url{https://www.sharelatex.com/learn/Bibtex_bibliography_styles}.

\section{Joitain ohjeita Latexilla kirjoittamiseen}

Seuraavassa on joitain ohjeita tämän tutkielmapohjan käyttöön maisterintutkielmassa. Kirjoittamisohjeita löytyy useasta eri lähteestä. Voit esimerkiksi tutustua kandidaatintutkielman ohjeisiin. 
Ohjaajan kanssa on hyvä keskustella aikaisessa vaiheessa työn rakenteesta.

Yksityiskohtaisia ohjeita Latexin käyttämäsestä saa parhaiten hakemalla verkosta, esim. haku englanniksi "Overleaf table positioning" tuottaa oletettavasti aika toimivan vastauksen.

\section{Kuvat}
Kuva~\ref{fig:logo-fi} toimii esimerkkinä kuvan lisäämisestä työhön (katso tarkemmin mallia Latex lähdekoodista). Muista myös viitata jokaiseen kuvaan tekstissä. 

\begin{figure}[ht] % remove [h!] for automatic placement, which is probably better for a thesis with more text on page
\centering 
\includegraphics[width=0.3\textwidth]{template/figures/HY-logo-ml.png}
\caption{Helsingin yliopiston logo matemaattis-luonnontieteellisen tiedekunnan värein.\label{fig:logo-fi}}
\end{figure}

\newpage % just to keep the table on the same page with the short piece of text
\section{Taulukot}

Taulukossa~\ref{table:results-fi} on esimerkki kokeellisten tulosten raportoinnista taulukkona. Muista myös viitata jokaiseen taulukkoon tekstissä.
\begin{table}[ht]
\centering
\caption{Kokeelliset tulokset.\label{table:results-fi}}
\begin{tabular}{l||l c r} 
Koe & 1 & 2 & 3 \\ 
\hline \hline 
$A$ & 2.5 & 4.7 & -11 \\
$B$ & 8.0 & -3.7 & 12.6 \\
$A+B$ & 10.5 & 1.0 & 1.6 \\
\hline
%
\end{tabular}
\end{table}



% BSc instructions
%\chapter{Johdanto}


Kaikessa julkaistavaksi tarkoitetussa tekstissä kirjoittajan luomisen ja
esitystavan vapautta rajoittavat monet ohjeet ja tarkatkin määräykset.

Parhaimmillaan lukijalle ja kirjoittajalle yhteinen, tuttu säännöstö luo
eräänlaisen tukiverkoston, joka tukee sanoman siirtymistä vääristymättä.
Kirjoituksen lukija löytää kirjoituksesta helpommin olennaisen sisällön,
jos kirjoituksen ulkoasu ja sisällön rakenne vastaavat hänen
tottumuksiaan. Sama koskee myös kirjoittajaa. Noudattaessaan valmista
esitystapamallia kirjoittajan ei tarvitse käyttää aikaansa itse työn
kannalta toissijaisten seikkojen miettimiseen, vaan hän voi keskittyä
hiomaan tekstin sisältöä. Siksi kannattaa harjoitella myös työn ulkoasua
koskevien ohjeiden noudattamista, vaikka omasta mielestään osaisikin
valita esitykselleen ohjetta paremman muodon.

Tämä kirjoitus on tarkoitettu Helsingin yliopiston
Tietojenkäsittelytieteen osastoon alempien opinnäytteiden ja
harjoitusten ulkoasun ja rakenteen ohjeeksi. Ohje soveltuu siten
kandidaatintutkielman kirjoittamisen kurssille, ohjelmistotuotantoprojekteihin, seminaareihin ja
pro gradu -tutkielmiin. (Kirjoitus on päivitetty uusintapainos aiemmista
ohjeista, jotka kurssin luennoijat ovat laatineet \citep{erkio01,erkiomakela96,erkio94,verkamo92}.)

Tyylimäärittely on saatavissa pdflatex- ja word-versiona.
Tyylimäärittelyitä valitessa on huomattava ohjeet tekstien syöttöön
liittyvässtä koodauksesta (UTF8,ISO 8859-15).
Tämän kirjoituksen tukena sopivat käytettäväksi tavanomaiset latex- tai
word-oppaat. 

\chapter{Kirjoituksen rakenne}

Tarkastellaan aluksi tieteelliseltä tekstiltä odotettuja
kirjoituksen osia. Samoihin asioihin on luonnollisesti syytä
kiinnittää huomiota myös muussa teknisessä kirjoittamisessa. Huomattakoon, että tämä teksti ei ole tieteellinen teksti, eikä siten itse sisällä kaikkia niitä elementtejä, jotka tieteellisen tekstin sisällölliseen antiin kuuluvat. Tällaisia puutteita ovat esimerkiksi johdannon tutkimuskysymyksen asettelun puuttuminen sekä arvoivan materiaalin puute tekstin lopussa, sekä yhteenvedon latteus.
Teksti rajoittuu siten otsikkonsa mukaisesti vain tekniseen sisällön asetteluun.

\section{Tiivistelmä}



Tiivistelmäsivu sisältää seuraavat osat: työn bibliografiset tiedot,
tiivistelmäteksti, aiheluokat ja avainsanat. Bibliografiset tiedot
koostuvat työn otsikosta, tekijän nimestä, julkaisupaikan tiedoista,
julkaisuajankohdasta ja sivumäärästä.

Tiivistelmäteksti on lyhyt, yleensä yhden kappaleen mittainen
(maksimissaan noin 100 sanaa) selvitys
kirjoituksen tärkeimmästä sisällöstä: mitä on tutkittu, miten on
tutkittu ja mitä tuloksia on saatu.


Aiheluokat kuvataan ACM Computing Classification System -luokituksen (CCS)
luokituksen mukaisesti. Luokittelussa käytetään täysia polkuja juurisolmun CCS osoittamista lähtöposteistä lehtisolmuihin. Polkuja voi antaa 1-3 aihepiirien soveltuvuuden mukaan, mitä alempi opinnäyte, sen vähemmän polkuja se tarvitsee. 
Poluissa tasot erotetaan toisistaan nuolella eteenpäin. Kun  polun nimisanoja arvioidaan suhteessa työn sisältöön, merkitään boldface-fontilla tärkein termi, italics-fontilla toiseksi tärkein. Näin menetellään, mikäli jotkin termeistä ovat olennaisesti paremmin kuvaavia kuin muut polun termit. Nimettyjen polkujen lisäksi lukija voi siten tarkastella lisäulottuvuutena myös tärkeiksi merkittyjen termien joukkoa sinänsä.
Avainsanoiksi valitaan kirjoituksen sisältöä
keskeisesti kuvaavia käsitteitä.

\section{Johdanto}


Johdannon tarkoituksena on kertoa yleiskielisesti
työn tavoite. Kerrotaan (kuten tiivistelmässäkin, mutta laveammin),
mitä on tutkittu, miten on tutkittu ja mitä tuloksia on saatu.
Jotta kysymyksenasettelu ja tulokset on lukijan helppo oikein tulkita on syytä aloittaa johdanto asettelemalla tutkimus asiayhteyteensä, esimerkiksi kertomalla aluksi, minkälaisessa yhteydessä tarkasteluun otettavat haasteet esiintyvät ja keiden on ratkaisuista tarkoitus hyötyä.

Johdannon pituus määräytyy suhteessa koko kirjoitelman pituuteen.
Parisivuinen kirjoitus ei erikseen otsikoitua johdantoa kaipaa, sillä
se itsessään on laajennettu tiivistelmä. Kymmensivuisen
kirjoituksen johdanto voi olla vaikkapa sivun tai puolentoista
mittainen. Pro gradu -tutkielman 50-70-sivuiseen kokonaisuuteen
tuntuu 2-4-sivuinen johdanto kohtuulliselta. 

Johdanto kertoo siis lyhyessä, yleistajuisessa muodossa
koko kirjoitelman kysymyksenasettelun, juonen sekä tulokset ja johtopäätelmät.
Tämän luettuaan lukija voi päätellä, haluaako syventyä asiaan tarkemmin
lukemalla koko kirjoituksen.


\section{Käsittelyluvut}

Käsittelylukujen työnjako määräytyy käsiteltävän asian luonteen
mukaisesti.
Lukijan ohjailemiseksi kukin pääluku kannattaa aloittaa lyhyellä
kappaleella, joka paljastaa mikä kyseisen luvun keskeisin sisältö on ja
kuinka aliluvuissa asiaa kehitellään eteenpäin.
Erityisesti kannattaa kiinnittää huomiota siihen, että lukijalle ilmaistaan selkeästi miksi kutakin asiaa käsitellään ja miten käsiteltävät asiat suhtautuvat toisiinsa. 

Jäsentelyongelmista kielivät tilanteet, joissa
alilukuja on vain yksi, tai joissa käytetään useampaa kuin
kahta tasoa (pääluku ja sen aliluvut). Kolmitasoisia
otsikointeja saatetaan tarvita joissakin teknisissä
dokumenteissa perustellusti, mutta nämä muodostavat poikkeuksen.

Perusohjeena on käyttää tekstin rakenteellisesti painokkaita paikkoja,
kuten lukujen avauksia ja teksikappaleiden aloitusvirkkeitä
juonenkuljetukseen ja informaatioaskeleiden sitomiseen toisiinsa.
Tekstikappaleiden keskiosat, samoin kuin lukujen keskiosat selostavat
asiaa vähemmän tuntevalle yksityiskohtia, kun taas aihepiirissä jo
sisällä olevat lukijat voivat alkuvirkkeitä silmäilemällä edetä
tekstissä tehokkaasti eksymättä tarinan juonesta.

Kullakin kirjoittajalla on oma temponsa, joka välittyy lukijalle
tekstikappaleiden pituudessa ja niihin sisällytettyjen ajatuskulkujen
mutkikkuudessa. Kussakin tekstikappaleessa pitäisi pitäytyä vain yhdessä
informaatioaskelessa tai olennaisessa päättelyaskelessa, muuten juonen
seuraaminen käy raskaaksi olennaisten lauseiden etsiskelyksi. Yksivirkkeisiä
tekstikappaleita on syytä varoa. 



\section{Lähdeviittausten käyttö}


Olennaisia opittavia asioita viittaustekniikoissa ovat viitteen paikka
tekstissä, oikea lähdeluettelojärjestys valitun viitetyylin parina sekä
taito ja tahto noudattaa annettua tyylimääräystä. Väitöskirjoissa ja
lehti- tai konferenssiartikkeleissa tekstin hyväksyminen riippuu myös
näiden yksityiskohtien asianmukaisesta käsittelystä. Tästä syystä
laitoksella nähdään tarpeelliseksi opiskelijoiden tutustua edes
pinnallisesti myös muihin tyylilajeihin ja oppia käyttämään
automatisoituja muotoilutyökaluja tehokkaasti, jolloin tyylimuutokset
ovat tehokkaita.

Lähdeviitteet sijoitetaan aina virkkeen sisäpuolelle. Siten esimerkiksi
tekstikappaleen lopussa irrallaan oleva viite ei ole asiallinen. Tilanne
ei muutu, vaikka viite sujautettaisiin tekstikappaleen viimeisen
virkkeen sisään. 
Lähdeviittauksen yhteyteen merkitään mukaan tarkentavat
sivunumerot, mikäli lukijan olisi työlästä löytää asianomainen kohta
viitatusta lähteestä. 


Tehokkaita viitteensijoittelupaikkoja ovat esimerkiksi uuden käsitteen
nimeämiskohta ja virkkeen loppu kun kyseessä on lähteestä lainattu
väite. On myös muistettava lainausmerkkien käyttö silloin kun tehdään
suoria lainauksia.

Tekstin jäsentelyn on tuotava selkeästi esiin, mihin asiaan viite
liittyy. Samalla tulee ymmärrettäväksi se, kuinka pitkään
tekstikatkelmaan ko. viite liitetään. Ei ole siten asiallista aloittaa
lukua nimeämällä yhtä tai useampaa lähdettä luvun taustaksi, vaan
viitteitä on kiinnitettävä täsmällisemmin väitteisiin ja käsitteisiin.
Luvun avaus viitetiedolla voi olla oire myös suuremmasta ongelmasta:
lähderiippuvuudesta. Aloitteleva kirjoittaja helposti toistaa lähteestä
oppimaansa ilman että tarpeellinen analysointi ja prosessointi suhteessa
muuhun opittuun olisi vielä tapahtut.

Viitteillä ja sanamuodolilla on
myös tuotava selkeästi esiin se, mikä teksissä on lainattua ja mikä oman
pohdinnan ja valikoinnin tulosta.


Lähdeviittauksiin käytetään Tietojenkäsittelytieteen osastolla
numeroitua tyyliä ja APA-tyyliä, valinnan näiden välillä tekevät kunkin
ryhmän valvoja ja ohjaaja yhdessä. 
Numeroitu tyyli on esimerkiksi IEEE- ja
ACM-julkaisuissa yleisesti käytetty ja puolustaa siten paikkaansa.
APA-tyyli on poikkeuksellinen ns. kovissa tieteissä, mutta monet
valvojista pitävät siitä sen luettavuuden vuoksi. Numeroita joutuu nimiä
useammin tarkistamaan lähdeluettelosta, sillä tarkastus- ja
arvointiprosessiin kuuluu arvioida myös lähteiden valitaa ja niiden
käyttötapaa.





\section{Yhteenveto}

Yhteenveto  vaatimattomimmillaan on vain lyhyt kertaus kirjoituksen
keskeisistä asioista. Arvokkaamman yhteenvedon saa aikaan kommentoimalla
 työn tulosten arvoa, työn liittymistä ympäristöön ja
tulevaisuudennäkymiä. Tällaiset arviot  huolellisesti
perusteltava.

\section{Lähdeluettelon laatiminen}

Tieteellisen kirjoittamisen kurssin töiden lähdeluetteloiden
laatimisessa noudatetaan seuraavia ohjeita.

Niiden taustalla on kaksi
keskeistä pyrkimystä: tehdä viitatun lähteen hankkiminen luettavaksi
mahdollisimman helpoksi ja ilmaista, millaisen arviointiprosessin
läpi käyneeseen kirjoitukseen vedotaan.
Näistä syistä
\begin{itemize}
\item lähdeviitteen tulee aina olla niin tarkka, että
lähde on sen perusteella tunnistettavissa ja löydettävissä luetteloista
ja kirjastoista,
\item erityyppisten lähteiden (kirjat, konferenssit, lehdet) on erotuttava
toisistaan
ja
\item luettelon eri osien tulee olla mahdollisimman
yhdenmukaisia, erityisesti lähdetyypin sisällä.
\end{itemize}


Riippumatta käytettävästä viitetyylistä, 
lähteet ovat Tietojenkäsittelytieteen osaston opinnäytteiden lähdeluetteloissa tekijän nimen mukaisessa aakkosjärjestyksessä,
saman tekijän (tekijäryhmän) työt julkaisuajan mukaisessa
järjestyksessä. Jos jollakin lähteellä ei ole henkilötekijää, se
aakkostetaan julkaisun nimen mukaisesti. 

Kustakin lähteestä annetaan seuraavat tiedot, edelleen viitetyylistä riippumatta:
\begin{itemize}
\item (tarvittaessa lähdeviitelyhenne).
\item
tekijän tai tekijöiden nimet (sukunimi, etunimien alkukirjaimet)
alkuperäisessä järjestyksessään; jos tekijöitä on enemmän kuin kolme,
voidaan toimia siten, että
vain ensimmäinen tekijä nimetään ja muiden tilalle kirjoitetaan {\em et
al.}
\item
julkaisun tai artikkelin nimi alkuperäisessä muodossaan
\item
julkaisupaikan tiedot:
\begin{itemize}
\item
kirjasta: kustantaja, julkaisupaikka (voidaan jättää pois, jos kyseessä
on tunnettu kustantaja), vuosi ja
\item
lehtiartikkelista: lehden nimi, volyymi, numero, vuosiluku ja kuukausi (suluissa),

\item
artikkelikokoelmassa (esim. konferenssijulkaisussa) ilmestyneestä
artikkelista:
\begin{itemize}
\item kokoelman nimi, toimittaja, kustantaja, julkaisupaikka ja vuosi
{\em tai}
\item konferenssin nimi, järjestäjä, paikka ja aika,
\end{itemize}
\item
raportista: julkaisusarja, raportin numero, julkaisupaikka, julkaisija ja vuosi
ja
\item
www-lähteestä: verkko-osoite, voimassaoloajankohta, mahdollisesti
viittausajankohta hakasuluissa
\end{itemize}
\item
sivunumerot, mikäli lähteenä käytetty julkaisu on artikkeli tai kokoomateoksen itsenäinen luku.
\end{itemize}

Normaaliin suomalaiseen tapaan artikkelin nimessä ainoastaan
ensimmäinen sana kirjoitetaan isolla alkukirjaimella, sen sijaan
konferenssien ja kokoelmajulkaisujen nimissä käytetään isoa
alkukirjainta jokaisen sanan alussa (artikkelisanoja ja prepositioita
lukuunottamatta). Katso mallia oheisista esimerkeistä.
Kokoelman nimen edessä on syytä selvyyden vuoksi käyttää sanaa {\em
Teoksessa}, paitsi kun on kysymys konferenssijulkaisusta, jonka nimi
alkaa lyhenteellä {\em Proc.} (sanasta Proceedings). Tällöin ei tarvita
mitään täydennystä.
Tämän eron näkee esimerkiksi vertaamalla
lähdeviitteiden~''\citep{dantowsley90}''
ja~''\citep{gannonetal89}'' ulkoasuja.

WWW-lähteiden käytössä on syytä muistaa, että verkossa julkaisukynnys on
olematon. Kannattaa siten keskittyä tunnettujen tieteellisten
kustantajien julkaisuihin ja niihin teknisiin standardeihin, joille WWW
on ainoa julkaisukanava. Mikäli sama julkaisu on saatavissa myös
perinteisessä muodossa, viitataan ensisijaisesti siihen ja käytetään
verkko-osoitetta lisätietona. Lähdeluettelossa on annettu esimerkit
useita kanavia julkaistusta kirjoituksesta~\citep{abiteboul,dietinger} sekä pelkästään
WWW-julkaisuna
leviävästä standardista~\citep{bray}.

Erityisesti varoitetaan Wikipedian käytöstä tieteellisessä tekstissä.
Vaikka sen avulla on helppo alustavasti tutustua joihin aihepiireihin ja
asiantuteva lukija voisi teksin kelvolliseksi tiettynä hetkenä
hyväksyäkin, ei se foorumina millään lailla täytä tieteellisesti
vertaisarvoidun tutkimusfoorumin kriteerejä.
Jos Wikipedia-artikkelia  ei mitenkään malta ajankuvana olla mainitsematta, käytettäköön jotain muuta kuin lähdeviitetekniikkaa tähän taiteelliseen otteeseen, vaikkapa alaviitteitä. Olennaista silloinkin on, että tieteellinen sisältö ei tule tällä korvatuksi vaan sen puute korostetuksi.

WWW-lähteeseen viittaamisessa pätevät samat periaatteet kuin
perinteisiin lähteisiin viitattaessa: lähdeviitteessä ilmaistaan
otsakkeet, kirjoittajat, toimittajat js muut seikat. Eroa on ainoastaan
verkko-osoitteen ja sen voimassaoloajankohdan ilmaisemisessa. Mikäli
lähde on julkaistu ainoastaan verkossa, voidaan web-osoitetta (URL)
käyttää vastaavasti kuin perinteisen julkaisun paikannusinformaatiota
(lehden ja se numeron julkaisutiedot). Lähdeluettelossa on WWW-viittausten yhteydessä aina syytä ilmaista päivämäärä, jolloin linkin voimassaolo ja lähteen sisältö on tarkastettu.
Esimerkkeinä verkkoviitteistä soveltuvat seuraavat:
\begin{itemize}
\item Gergen, Kenneth (1999) Narrative, Moral Identity and Historical
Consciousness: a Social Constructionist Account.
http://www.swarthmore.edu/SocSci/kgergen1/text3.html. Haettu 11.6.1999.
\item 	
Ritala-Koskinen, Aino and Valokivi, Heli (2006) The Role of Development
Skills in Social Work Practice Education in Finland. Social Work and
Society, The International Online-Only Journal 4(2006)1.
http://www.socwork.net/2006/1/series/transition/ritalakoskinenvalokivi.
Viitattu 30.8.2006.
\item Heinisuo, Rami and Ekholm, Kai (1997) Elektronisen viittaamisen
opas. Jyväskylän yliopiston kirjaston julkaisuja n:o 40. Jyväskylä:
Jyväskylän yliopiston kirjasto. http://www.pori.tut.fi/~multisil/evo/.
Viitattu 29.8.2006.
\end{itemize}

Kirjoituksen lähdeluettelossa luetellaan täsmälleen ne lähteet, joihin
viitataan kirjoituksen tekstiosassa. Tämän kirjoituksen lähdeluettelo on
tarkoitettu lähinnä esitystavan esimerkiksi, mistä syystä siinä on
''ylimääräisiä'' lähteitä.


Pääsääntöisesti julkaisun tai artikkelin nimen perään tulee piste,
samoin kunkin lähteen bibliografisten tietojen perään. Muut tiedot
erotetaan toisistaan pilkulla. Useimmissa tapauksissa  
voidaan noudattaa teknisten välineiden antamaa mallia, sillä edellytyksellä, että ylläolevat vaatimukset muuten täyttyvät.



\chapter{Ulkoasulliset seikat}

  Tässä luvussa käsitellään yleisimpiä
tekstin tekniseen esittämiseen liittyviä seikkoja.  

\section{Työn osien järjestys}

   Kirjoituksen alussa on aina
erillinen, mallin mukainen kansilehti. Toisena sivuna on
tiivistelmäsivu, sen jälkeen sisällysluettelo (yksi tai useampia
sivuja) ja sitten varsinainen teksti. Sivunumerointi aloitetaan vasta
ensimmäiseltä tekstisivulta (arabialaisella ykkösellä). (Tarkat
jättävät ykkössivun numeromerkittä.) Sisällysluetteloon merkitään
kaikki (numeroidut) otsikot ja vastaavat sivunumerot. Monet
tekstin\-käsittelyjärjestelmät muodostavat itse sisällysluettelon,
jolloin kirjoittajan ei tarvitse huolehtia luettelon sivunumeroiden
päivittämisestä tekstin kehittyessä. Sisällysluettelo\-sivu ja sitä
edeltävät sivut voidaan haluttaessa numeroida erikseen (roomalaisin
numeroin) esimerkiksi tämän mallin mukaisesti.  

Varsinaisen tekstin
jäljessä, mutta itse työhön kuuluvana on ensimmäisenä lähdeluettelo,
jonka otsikkoa ei numeroida. Lähdeluettelon jälkeen sijoitetaan
mahdolliset liitteet, jotka otsikoidaan ja varustetaan sisäisillä
sivunumeroilla.  



Mikäli kuvista, algoritmeista
ja taulukoista halutaan tehdä yhtenäinen luettelo, sijoitetaan
luettelot sisällysluettelon jälkeen. Luetteloiden käyttöarvosta on
eriäviä mielipiteitä, joten niiden laatimiseen ei varsinkaan ilman
tekstinkäsittelyjärjestelmän tukea kannata ryhtyä ilman tarkastajan
erityistä toivetta.  

Mikäli kirjoitukseen erityissyistä halutaan
liittää aakkosellinen hakemisto, sijoitetaan se lähdeluettelon jälkeen
ennen liitteitä. Indeksi merkitään sisällysluetteloon samoin kuin
lähdeluettelo (numeroimaton luku). Mikäli indeksin tekemiseen
ryhdytään, on syytä käyttää tekstinkäsittelyjärjestelmän tarjoamaa
automatiikkaa.

Teksin luonnollisen juonenkuljetuksen mukana esiin
tulevien käsitteiden määrittelyjen sijasta ei pidä yrittää sen enempää
pakata kaikkia määritelmiä johdantoon kuin laatia johdantoa ennen
käsitelistaa tai lyhenteiden selityslistaa. Kumpikaan ei sovi
tavanomaiseen argumentoivaan tieteelliseen tekstityyliin, vaikka
teknisessä yhteydessä niillä liitteinä voi olla lisäarvoa.


\section{Tekstin yleinen sijoittelu}

Lopullinen tutkielmaversio voi olla yksi- tai kaksipuoleiseksi aseteltua
ja riviväliltään 1,5 tai 1.  Erityyppisissä
teksteissä haasteet ja asetteluvaatimukset voivat olla erilaiset. Erota
kappaleet toisistaan yhdellä tyhjällä rivillä tai
käytä tekstinkäsittelytyökalujen ominaisuuksia
hyödyksesi ja määrittele tekstikappaleiden väliin jäävä tila hieman
normaalia riviväliä suuremmaksi.

Kirjoituksen lukujen, kuvien ja taulukoiden erottumisen kannalta
tärkein keino on riittävän tilan käyttö niiden ympärillä. Kuvan ja
nimekkeen tulee olla selkeästi yksi kokonaisuus, joka eroaa muusta
tyhjän tilan rajaamana. Kuvan tai taulukon on aina numerointinsa ja
nimekkeensä kanssa mahduttava yhdelle sivulle tai varmasti
kaksipuolisena paperidokumenttina tarkasteltavassa tekstissä aukeamalle.
Kuvissa fonttikoko ei saa alittaa 8 pistettä.



Jos uusi luku tulisi alkamaan aivan sivun alareunasta (vain yksi tai
kaksi riviä varsinaista tekstiä), aloita mieluummin uusi sivu. Jokaista
uutta lukua ei kuitenkaan ole tarpeen --- etenkään lyhyessä
kirjoituksessa --- aloittaa uudelta sivulta: jos kirjoituksessa on
paljon melkein tyhjiä sivuja, lukija voi epäillä, että kirjoittaja on
yrittänyt saada kirjoituksensa näyttämään pitemmältä kuin se onkaan. 
 Tyhjää tilaa kannattaa käyttää hyödyksi myös kuvien ja taulukoiden
yhteydessä. Erityisesti jos kirjoituksessa käytetään kauttaaltaan samaa
tekstityyppiä, tyhjät rivit ovat välttämättömiä erottamaan esimerkiksi
tekstiä ja taulukkoa toisistaan. Tyhjä tila on halpaa, mutta se lisää
selkeyttä ja luettavuutta.  


\section{Kuvat ja taulukot}


Kuva tai taulukko sijoitetaan mahdollisimman lähelle
(ensimmäistä) tekstikohtaa, jossa siihen viitataan, ei kuitenkaan
kyseistä viittausta aikaisemmaksi.
Tekstissä on syytä myös kertoa, mitä kuvalla halutaan havainnollistaa.
Kuvan voi lukea monella eri tavalla, joten lukijaa on ohjattava.

Kuvaa ei pidä sijoittaa välittömästi luvun otsikon alle, vaan on
aloitettava tekstillä. Kuvaa ei pidä sijoittaa keskelle tekstikappaletta
(saati virkettä), paitsi jos kuva tulee sivun alkuun tai loppuun eikä
kappaleen jatkumisesta tule epäselvyyttä.

Kuvan ei aina tarvitse olla välittömästi viittaavan kappaleen
perässä. Esimerkiksi viittauskohdan ja
vasta seuraavalle sivulle mahtuvan kuvan väliin jäävää sivun loppuosaa
ei jätetä tyhjäksi. Kuvaa ei kuitenkaan pidä viedä seuraavaa
sivua kauemmas viittauskohdasta.


Varsinaista kuvan esittämistä havainnollistaa kuva~\ref{kuvaesimerkki}.
Huomiota on kiinnitettävä kuvan osien ja tekstimerkintöjen näkyvyyteen,
kuvan numerointiin ja otsikointiin. 

\begin{figure}[ht]
%\begin{figure}[tbh] t= top, b = bottom, h=here
\ \newline
\begin{center}
\includegraphics[width=0.75\textwidth]{kuvaesimerkki.pdf}
\caption{Kuvan elementit.}
\label{kuvaesimerkki}
\end{center}
\end{figure}


Kuvien kokoon on kiinnitettävä huomiota. Käytettyjen merkintöjen
on oltava helposti luettavissa ja selkeät. Esimerkiksi
suorituskykykäyriä esitettäessä akselit on nimettävä, asteikot
merkittävä ja käytetyt yksiköt tuotava selkeästi esiin.
Samankaltaisia asioita esitettäessä useammalla kuvalla on
syytä käyttää samaa mittakaavaa vertailun helpottamiseksi.

Kuvan otsikko kirjoitetaan kuvan alle ja sen tulee olla mieluummin lyhyt
ja ytimekäs kuin liian selittelevä.
Samoin toimitaan taulukoiden otsikoinnissa.

Kuvat ja taulukot numeroidaan juoksevasti. Pitkissä teksteissä käytetään
kaksitasoista numerointia (esimerkiksi Kuva 3.1) pääluvuittain, lyhyissä
riittää yksitasoinen numerointi.


Kuva- ja taulukko-otsikoiden yhdenmukaiseen esitystyyliin on syytä kiinnittää
huomiota, samoin mm. välimerkkeihin. Luontevaa on käyttää
kuvatekstin lopussa pistettä, ovathan useimmat kuvateksteistä virkkeitä. 

(Kuvien ja taulukoiden otsikointityyli vaihtelee
kustantajittain ja julkaisuittain. Samoin tuntuu suositeltava käytäntö
Tietojenkäsittelytieteen laitoksen sisällä vaihtelevan taulukon
otsikon sijainnin suhteen.)


\section{Otsikot}

Otsikoissa voi käyttää muusta tekstistä poikkeavaa kirjasintyyppiä,
alleviivausta, suurempaa kirjasinkokoa tms.\ erotuskeinoa, yleensä
kuitenkin vain yhtä näistä, koska kovin monta erilaista kirjasintyyppiä
ja -kokoa tekee ulkoasusta helposti sekavan.  Otsikoiden esitystavan on
oltava johdonmukainen läpi koko kirjoituksen. Numeroimattomia
''ylimääräisiä'' otsikoita ei tule yleensä käyttää.


\section{Mallin käyttö}

Voit käyttää tätä kirjoitusta mallina oman opinnäytteesi ulkoasua
varten. Eri tekstinkäsittelyjärjestelmissä käytössä olevat yksityiskohdat kuten
kirjasintyypit ja -koot ja rivivälit  poikkeavat toisistaan, joten
pienet poikkeamat ovat toki hyväksyttäviä.

Tieteellisen kirjoittamisen kurssin luennoilla ja
liitteenä olevassa ohjeessa annetut töiden ohjeelliset sivumäärät
koskevat työtä, joka vastaa ulkoasultaan tätä ohjetta (kirjasinkoko
12~pistettä). Tässä tekstissä keskimääräinen rivin pituus lienee noin
80~merkkiä ja sivun pituus 35-40~riviä.
Sivumääriin lasketaan varsinaisen tekstiosuuden pituus ja lähdeluettelo
(arabialaisin numeroin numeroitu osuus), ei kansilehteä, tiivistelmää
eikä sisällysluetteloa. Sivumääräarviossa otetaan huomioon hyvin vajaat
sivut, joita syntyy paljon lyhyiden lukujen ja taittotyyliin määritellyn
luvun avauksen pakottaminen oikeanpuolimmaiselle sivulle. 

\chapter{Yhteenveto}

Tämän kirjoituksen tarkoituksena on toimia muistilistana eräistä
esitystavallisista säännöistä, joihin harjoitusten ja tutkielmien
kohdalla on syytä kiinnittää huomiota.

Annetut ohjeet on laitoksen henkilökunta muotoillut yhdessä keskustellen
ja noudattaen oman tieteenalansa perinteitä. Eri erikoistumisaloilla ja
erilaisillaa määräävässä asemassa olevissa julkaisufooruilla käytänteet
vaihtelevat ja nuorten tutkijoiden onkin tiedostettava ero yleisten
sisältöohjeiden ja teknisten muotoilusääntöjen välillä. Aina tekstin
valmistuessa on tarkastettava erikseen, täyttääkö se annetut
pituusrajoitteet ja vastaako se annettuja muotoiluohjeita, olivatpa ne
kuinka pikkutarkkoja tahansa. Tarkasta sääntöjen noudattamisesta syntyy
yhteinäisyyttä kokoovan julkaisun tasolla, mikä helpottaa lukijoiden
työskentelyä.

Tämä ohje vastaa vain asettelullisiin kysymyksiin ja sen rinnalla on
syytä tutustua materiaaliin ja luentoihin, joissa keskitytään tekstin
varsinaiseen sisältöön. Olennaisin väline on kuitenkin akateemisesti
pidemmälle ehtineen, jo julkaisuja rakentaneen ohjaajan palaute ja
mentorointi.


%
\chapter{Introduction}


In all writing for publication, the writer's freedom of creation and expression are limited by a number of guidelines and specific regulations.

At best, a familiar set of regulations shared by reader and writer can create a kind of support network that allows the message to be relayed without distortion. It will be easier for readers to find 
the pertinent contents in a piece of writing if its layout and structure are the same as they are used to. This also applies to writers. When writers follow a set presentation model, 
they do not have to waste time on considerations that are secondary to the work itself, but they can concentrate on polishing the contents of the text. This means that it is a good 
idea to practice following the rules for the layout, though you may think you know how to select a better way to present your work.

This is a guide for the layout and structure of theses and essays at the Department of Computer Science at the University of Helsinki. It is thus applicable to the course 
Scientific Writing, the software engineering projects, seminars, and MSc theses. (This is an updated version of 
the previous guide written by the course lecturers \citep{erkio01, erkiomakela96, erkio94, verkamo92}.)

The \LaTeX\ guide and \LaTeX\ style that has been  published on the department's web site can be used as support for this guide.


\chapter{Structure}

Let us start by looking at the sections expected to be in a scientific text. Keep in mind that the same expectations go for all kinds of technical writing.
However, this document in itself, is not a scientific or a research
text, so there will be content lacks in terms of research question
setting in the introduction and evaluative material in the last sections.

\section{Abstract or summary}
%\enlargethispage{5mm}


The summary page contains the following elements: the bibliographical data of the work, an abstract, topic classification, 
and the key words. The bibliographical data consists of title, name of the author, place of publication, date of publication, and number of pages.

The abstract should be short, generally one paragraph (100 words maximum) explaining the main contents of the work: topic, methodology and results.


 Topics are classified according to the ACM Computing Classification System
(CCS). A small set of paths (1-3) should be used, starting from any top nodes
referred to bu the root term CCS leading to the leaf nodes. The elements
in the path are separated by right arrow, and emphasis of each element individually can be indicated
by the use of bold face for high importance or italics for intermediate
level. The combination of individual boldface terms may give the reader
additional insight. 

\section{Introduction}


The purpose of the introduction is to introduce the goals of the work in general terms. Describe topic, 
methodology and results (as in the abstract, but expand it).
In order to provide the reader a good starting point for
interpretations, it is good to start the introduction by
contextualisation of the challenges and solutions to be discussed. For
example, why a certain domain has a particular challenge and who are
intended to benefit from the solutions proposed.

The length of the introduction depends on the length of the whole work. A few pages of text does not need a 
separate introduction, since it is an expanded summary in itself. The introduction to a 10-page text can 
be 1--1.5 pages long. For a 50--70-page MSc thesis, a 2--4-page introduction seems reasonable.

The introduction should shortly describe the problem field of the whole work, the plot, and the conclusions, 
in general terms. After reading it, the reader may decide whether to go deeper into the topic by reading the whole text.


\section{Topic chapters}

The nature of the matter at hand determines how the topic chapters are disposed.
In order to guide the reader, it is a good idea to start each main chapter with a short paragraph 
on what the main topic of the chapter is and how it progresses from one sub-chapter to the
next. Especially relevant is to express how concepts, challenges,
solutions and research steps are bound together. There should be enough
guidance for the reader to allow expectation of the right storyline.


Basic rule for easy to follow text is to use the natural emphasis of the
text structure to support the content matter key concepts and thought
processes. This means using openings of sectiosn and text paragraphs for
key arguments and information moves, while the internal parts of
paragraps are filled with supporting aspects that those less familiar
with the topic area need. Those with better background knoweldge can
quickly skim trhough the text without loosing any essential arguments.

Each author has his or her personal ryhtm in the text, which is visible
for the reader as the length of text paragraphs and complexity of
thought chains within. A good policy is to take only one information
move or transition per paragraph. This way the text stays easy to follow.

Signs of problems with the disposition of the text are easily seen in texts
with only one sub-chapter, or with more than two chapter levels (main and sub-chapters). There may be justifiable reasons to use three-level 
headings in some technical documents. Single sentence text paragraphs
are also to be avoided.



%\pagebreak
\section{Reference usage}


Relevant learning targets include superficial knowledge of several
citation styles and capability (and willingness) to follow a given
style and ordering of entires and bibliographical details in the list of references.
These aspects are essential as the approval of a PhD manuscript or
journal article may depend on them.

Disregard of which style you use, 
references are always placed inside sentences.  This means that e.g. a separate reference at the end of a paragraph would be inappropriate.

The structure of the text must clearly show what the reference relates to.  At the same time, it 
shows how long a piece of the text that the reference relates to.

Efficient positions for citations are right after the introduction
(definition) of a concept, a methods or such, or in the end of a claim
from the reference material. Furthermore, if quoting verbatime one must
use citation marks.

The text structures and wordings, in addition to the location of
ciations must clearly express whether claims or arguments are 
authors' own, or if they come from the contribution of others. Thus is
is of bad style to open a section by listing the citations on which the
section is based on. Such method can further indicate more serious
problems like following reference material as it is instead of analysing
and synthetising material into new though processes. 


\section{Conclusion}

At its simplest, a conclusion is merely a weak revision of the main points of the text.  All more valuable 
conclusion sections contain comments on e.g. the value of results, how the work relates to its environment, or 
future visions. This kind of evaluations should be well-grounded in fact, though, or the conclusion 
might inadvertently seem comical. 

\section{Creating the list of references}

The following guidelines should be followed when creating lists of reference for the 
assignments during the course Scientific Writing.

The guidelines are backed by two main goals: to make it as easy as possible to find the 
referenced source, and to show what kind of evaluation process the referenced work has undergone. For these reasons
\begin{itemize}

\item the reference notes should always be so exact that the source can be recognized and found in catalogues and libraries

\item different types of sources (monographs, conferences, journals) have to be easy to distinguish from each other, and 

\item the different parts of the list must conform to each other, especially for each source type.
\end{itemize}


Independent of the citation style in use, 
the sources are listed alphabetically according to the author's name, and works by the same 
author (group of authors) in the order in which they have been published. If some publication 
does not have an individual writer, it is alphabetized according to its name.


The following information should be given on each source independent of
the citation style:
\begin{itemize}
\item
Name(s) of author(s) (surname, initial letters of first names) in the original order; if there are more 
than three authors, we can write the name of the first author and {\em et al.} instead of the other names.

\item
the name of the publication or article in its original form

\item
place of publication:

\begin{itemize}

\item
of monographs: publisher, place of publication (can be omitted if the publisher is well known), year
\item
of journal articles: name of journal, volume, issue, year and month (in parenthesis)

\item
of articles from article collections (such as conference publications):
\begin{itemize}

\item name of collection, editor, publisher, place of publication and year 

{\em or}

\item conference name, coordinator, place and time,

\end{itemize}

\item
of a report: series, report number, place, publisher and year

\item
of a web source: URL, validity date, possibly the date when referenced in square brackets

\end{itemize}

\item
page numbers, if the source is an article or constitutes a chapter in a compilation.
\end{itemize}


When using web sources, you should keep in mind that the threshold for publication on the web 
is non-existent. It is better to concentrate on the publications of well-known scientific 
publishers and the technical standards for which the web is the only publication channel. If 
the same publication is available on paper, refer to that primarily and add the URL as a complement.


The list of references gives an example of a text that has been published through many 
channels~\citep{abiteboul,dietinger}and another example that shows a standard that has 
been disseminated only through the web~\citep{bray}. For web-based
references it is important to give both the publication date and the
time of reading and interpreting the material. The content is prone to
change, and dating your use, you protect your interpretations for
unnecessary accusations of being faulty, if it is known which version
you used.



The list of references for a text should list exactly the sources that the text refers to. 
The list of references for this text is an example of how to present sources, 
which is why it contains ''extra'' sources.


In any style,
put a full stop after the name of a publication or article, as well as after the bibliographical 
data of each reference. Separate the other pieces of information with a comma. As is normal in 
Finnish, only the first letter of the first word in the heading is capitalized, but in the titles of 
conferences and compilations, each major word is capitalized (not articles or prepositions). See the 
appended example for a model. For the sake of clarity, it is best to write {\em In the work} before 
the name of a compilation, except in the case of conference publications where the name starts with 
the abbreviation {\em Proc.} (for Proceedings). In such cases, no complement is necessary. You can 
see the difference by comparing the layouts of the references ~''\citep{dantowsley90}'' and~''\citep{gannonetal89}'' .
In case of other citation styles, it is likely that you can trust the
results given by the automated reference management tools.

\chapter{Layout}

This chapter discusses the main issues of the technical presentation of a text.


\section{Disposition of text sections}


Each text should include a separate cover page, as in this model. The second page contains an 
abstract, followed by a table of contents (on one or more pages), and then the main text.  
Pagination starts on the first page of the main text (with the Arabic numeral 1). (Rigorous writers 
leave out the number from the first page.) All the (numbered) headings and their page numbers should be 
written into the table of contents. Many word-processing systems create the table of contents automatically 
so that the writer does not have to worry about updating page numbers as the work progresses. If writers 
wish to, they can paginate the table of contents and previous pages separately (with Roman numerals), e.g. like this model.

After the main text, but as part of the text body, comes the list of references; its heading is 
not numbered. Any possible appendices are added after the list of references, with headings and internal page numbers.

%\enlargethispage{5mm}

If you want to make a coherent list of figures, algorithms and tables, this list should be placed i
mmediately after the table of contents. The value of such lists is a matter of opinion, so there 
is no need to create one --- especially if your word-processing software does not support it --- unless your instructor explicitly asks for it.

If for some special reason you want to add an alphabetical index top your work, it should be placed after 
the list of references and before appendices. An index should be noted in the table of contents, as s
hould the list of references (unnumbered chapter). Writers who want to create an index should use the automatic support their word-processing system offes.

\section{General text layout}

Nowadays you may plan your text to be printed on single side or double
side, with line spacing one, for the final copy.

For reviewing and feedback purposes, check with your readers for their
preferences. Most read and comment on paper and need some space for
marking. 
For drafts you can use wide margins too. 
Both margins should also be fairly broad (ca.~3~cm): the wider left margin is needed for binding the text and the
right one for the instructor's comments.  Leave enough space (2--3~cm) at the top and bottom, as well.

The most effective way to distinguish features like chapters, figures, etc,  is to have enough space 
in your text.  Separate paragraphs from each other with one and chapters with a few empty lines.  If a 
new chapter is about to start at the bottom of a page (with only one or two lines of text), it is better 
to start it on the next page.  However, it is not necessary to start every chapter on a new page, especially 
not with a short text; if the text contains many pages that are nearly empty, readers might suspect that the 
writer has tried to make it look longer than it is.


Any empty space may be utilized for displaying figures and tables.  Especially if a text is written with the 
same type of text throughout, empty lines are necessary for separating e.g.~text from tables.  Empty space 
is cheap, but adds to the clarity and readability of a text.

\section{Figures and tables}

%\enlargethispage{5mm}

All figures and tables should be placed as near the (first) place in the text that refers to them, but 
not before this reference. The text should also explain what the writer wants to illustrate with the 
figure. Figures can be interpreted in many different ways, so the reader needs guiding.

Figures should never be placed immediately under the heading of a chapter, but chapters should always 
start with text. Figures should not be placed in the middle of a paragraph (much less a sentence), except 
if the figure is placed at the top or bottom of a page and it is clear that the paragraph continues.

Figures do not always have to be placed immediately after the paragraph that refers to them. If there is 
not enough space for a figure on the same page as the paragraph that refers to it, the rest of the page 
should not remain empty, Though the figure is inserted on the following page. However, the figure should 
never be more than one page away from the reference.


The image~\ref{kuvaesimerkki} shows how to present a figure. You must pay attention to the visibility of 
figure parts and text, to the numbering of figures, and captions. 

\begin{figure}[ht]
%\begin{figure}[tbh] t= top, b = bottom, h=here
\ \newline
\begin{center}
\includegraphics[width=0.9\textwidth]{kuvaesimerkki.pdf}
%\rotatebox{90}{\includegraphics[scale=.75]{kuvaesimerkki.pdf}}
\caption{Figure elements.}
\label{kuvaesimerkki}
\end{center}
\end{figure}


We must also pay attention to the size of figures. Any annotations must be clear and easy to read When 
presenting performance curves, for example, the axels have to be named, the scale notated, and the units 
clearly presented. If you present many things with similar figures, you should use the same scale for easy comparison.

The caption of a figure should be written underneath it, and it is preferable that it be short and to 
the point than too explicit. The same goes for table captions.

Figures and tables should be numbered progressively. In long texts, two-level numbering should be 
used (e.g. Figure 3.1) according to the chapter number, but in shorter texts, one-level numbering is good enough.


You should pay attention to presenting figure and table captions consistently, as well as to punctuation 
marks. It is natural to put a full stop after a caption, since they are most often full sentences. 

(The style for figure and table captions vary according to publisher and publication. The recommendations 
at the Department of Computer Science also seems to vary with respect to where the caption should be placed.)


\section{Headings}

You can use a different font in headings than in the rest of the text, or underlining, larger fonts, or other 
methods of emphasis, but generally it is preferred that you only use one of these methods, because if there are 
very many font types and sizes, the layout looks messy.  The format of headings must be consistent throughout 
the text. You should not use any unnumbered ''extra'' headings.


\section{Using this model}

You can use this text as a model for the layout of your own assignment. The font types and sizes, 
line spacing, etc vary according to word-processing system, so small deviations from the rule are acceptable.

The directive number of pages for written assignments given during the lectures for the Scientific-writing 
course and in this guide are applicable to assignments with a layout like this guide's (font size 12 points). The 
average line in this text should consist of about 80~characters and one page about 30~lines. The number of pages 
includes the text itself and the list of references (the part paginated with Arabic numbers), not the 
cover page, summary, or table of contents.

\chapter{Conclusion}

This text is a checklist for some of the rules governing written presentations, which you should 
keep in mind when writing exercises and theses.

This collection of advise has been compiled by staff members as the
result of discussions on general good writing habbits in computer
science. The consensus is that all young researchers and academic degree
holders must be able to seek and follow instructions of layout and
structuring of texts depending on the contribution they are working on. 
This set of instructions aims at unifying the looks of theses and other
study reports from the department, but it is also a representative of
the instruction sets students will meet later.

This set of instructions only address the looks of the document. Other
instructions must be sought, read and gained experience with in order to
learn to select the suitable semantical contents for the scientific texts.



\end{appendices}
%%%%%%%%%%%%%%%%%%%%%%%%%%%%%%%%%%%%%%%%%%%%%%%%%%%%%%%%%

\end{document}
